%%%%%%%%%%%%%%%%%%%%%%%%%%%%%%%%%%%%%%%%%%%%%%%%%%%%%%%%%%%%%%%%%%%%%
%%                                                                 %%
%% Please do not use \input{...} to include other tex files.       %%
%% Submit your LaTeX manuscript as one .tex document.              %%
%%                                                                 %%
%% All additional figures and files should be attached             %%
%% separately and not embedded in the \TeX\ document itself.       %%
%%                                                                 %%
%%%%%%%%%%%%%%%%%%%%%%%%%%%%%%%%%%%%%%%%%%%%%%%%%%%%%%%%%%%%%%%%%%%%%


%%======================================================%%
%% to compile with pdflatex/xelatex use pdflatex option %%
%%======================================================%%


\documentclass[12pt]{article}% Math and Physical Sciences Reference Style %% Uncomment this for submission
%\documentclass[pdflatex,sn-mathphys]{sn-jnl-no-heading}% Math and Physical Sciences Reference Style 

%%\documentclass[referee,sn-basic]{sn-jnl}% referee option is meant for double line spacing

%%=======================================================%%
%% to print line numbers in the margin use lineno option %%
%%=======================================================%%

%%\documentclass[lineno,sn-basic]{sn-jnl}% Basic Springer Nature Reference Style/Chemistry Reference Style



%%%% Standard Packages
%%<additional latex packages if required can be included here>
\usepackage[a4paper]{geometry}
\geometry{
 paperheight = 297mm, % Tamaño A4: 297mm
 paperwidth = 210mm, % Tamaño A4: 210mm
 textheight = 247mm, % Valor máximo: 247mm
 textwidth =160mm, % Valor máximo: 160mm
 centering, % Lo habitual
 top = 25mm, % Lo habitual: 25mm
% includehead, % Incluir la cabecera en el área del texto
}
\setlength{\parindent}{7pt}
\usepackage{amssymb}               % para cargar las fuentes de la A.M.S.

\usepackage{amsmath}               % para cargar las fuentes de la A.M.S.
\usepackage{hyperref}   % Autoref

\usepackage{ dsfont }
\usepackage{enumitem}            %autoref enumerate
\usepackage{faktor}


\usepackage{color}

% Shortcuts
\newcommand\N{\mathbb{N}}
\newcommand\Real{\mathbb{R}}
\newcommand\Comp{\mathbb{C}}
\newcommand{\id}[1]{\mathds{1}_{#1}}
\newcommand{\mb}[1]{{\mathbf{#1}}}
\newcommand{\sch}[1]{{\mathcal{S}(\Real^#1)}}
\newcommand{\dsch}[1]{{\mathcal{S}'(\Real^#1)}}
\newcommand{\wickdel}[2]{:#1:\rvert_{\scaleto{#2}{3pt}}}
\newcommand{\braket}[1]{{\langle #1\rangle}}
\newcommand{\bra}[1]{\langle #1\lvert}
\newcommand{\ket}[1]{\rvert #1 \rangle}

%%%%

%\usepackage[normalem]{ulem}

%\jyear{2022}%
%=====================================================%
%% THEOREMS
%%====================================================%
\usepackage{aliascnt}   % second counter for autoref

%\theoremstyle{thmstyleone}%
%%\newtheorem{theorem}{Theorem}%  meant for continuous numbers
\newtheorem{theorem}{Theorem}[section]% meant for section wise numbers
%% optional argument [theorem] produces theorem numbering sequence instead of independent numbers for Proposition
\newaliascnt{proposition}{theorem}
\newtheorem{proposition}[proposition]{Proposition}% 
\newaliascnt{corollary}{theorem}
\newtheorem{corollary}[corollary]{Corollary}% 


%%\newtheorem{proposition}{Proposition}% to get separate numbers for theorem and proposition etc.

%\theoremstyle{thmstyletwo}%

\newtheorem{remark}{Remark}%

%\theoremstyle{thmstylethree}%
\newtheorem{definition}{Definition}%
\newtheorem{example}{Example}% 

\aliascntresetthe{proposition}
\aliascntresetthe{corollary}


\providecommand*\theoremautorefname{{\bf Theorem}}
\providecommand*\propositionautorefname{{\bf Proposition}}
\providecommand*\corollaryautorefname{{\bf Corollary}}
\providecommand*\exampleautorefname{{\it Example}}
\providecommand*\definitionautorefname{{\bf Definition}}

\raggedbottom %The \raggedbottom declaration makes all pages the height of the text on that page. No extra vertical space is added.
%%\unnumbered% uncomment this for unnumbered level heads



\usepackage{authblk}

\begin{document} 

\title{Geometric flavours of Quantum Field theory on a Cauchy hypersurface.\\ Part I:  Gaussian analysis and other Mathematical aspects  }

%%=============================================================%%
%% Prefix	-> \pfx{Dr}
%% GivenName	-> \fnm{Joergen W.}
%% Particle	-> \spfx{van der} -> surname prefix
%% FamilyName	-> \sur{Ploeg}
%% Suffix	-> \sfx{IV}
%% NatureName	-> \tanm{Poet Laureate} -> Title after name
%% Degrees	-> \dgr{MSc, PhD}
%% \author*[1,2]{\pfx{Dr} \fnm{Joergen W.} \spfx{van der} \sur{Ploeg} \sfx{IV} \tanm{Poet Laureate} 
%%                 \dgr{MSc, PhD}}\email{iauthor@gmail.com}
%%=============================================================%%


\author[1,2,3]{{José Luis} {Alonso}}
\author[1,2,3]{{Carlos} {Bouthelier-Madre}}
\author[1,2,3]{{Jesús} {Clemente-Gallardo}}
\author[1,3]{{David} {Martínez-Crespo}}

%\equalcont{These authors contributed equally to this work.}

\affil[1]{{Department of Theoretical Physics},{University of Zaragoza}, {{Facultad de Ciencias, Campus San Francisco},{Zaragoza},{50009},  {Spain}}}

\affil[2]{{Institute for Biocomputation and Physics of Complex Systems}, {University of Zaragoza}, {{Edificio I+D-Campus Río Ebro}, {Zaragoza},{50018},  {Spain}}}

\affil[3]{{Center for Astroparticles and High Energy Physics CAPA}, {University of Zaragoza}, {{Facultad de Ciencias, Campus San Francisco}, {Zaragoza}, {50009},  {Spain}}}

\maketitle



%%==================================%%
%% sample for unstructured abstract %%
%%==================================%%

\abstract{
    In this series of papers we aim to provide a mathematically comprehensive framework to the hamiltonian pictures of quantum field theory in curved spacetimes. Our final goal is to study the kinematics and the dynamics of the theory from the point of differential geometry in infinite dimensions.    
    In this first part we introduce the tools of gaussian analysis in infinte dimensional spaces of distributions that serve the basis to understand the Schrödinger and Holomorphic pictures, over arbitrary Cauchy hypersurfaces, using tools of Hida-Malliavin calculus. Here the Wiener-Ito decomposition theorem provides the  QFT particle interpretation. 
    Special emphasis is done in the applications to quantization of these tools in the second part of this paper. 
    We devote a section to introduce  Hida test functions as a notion of {\it second quantized } test functions. We also analyze of the ingredients of classical field theory modeled as distributions paving the way for quantization procedures that will be analyzed in \cite{alonso2023a}.  
}



\textbf{Keywords: }{Quantum fields,  Malliavin calculus, Schrödinger Picture, Holomorphic Picture, Momentum-field Picture, Infinite dimensional Fourier transform, Infinite dimensional Segal-Bargmann Transform, Infinite dimensional geometry}

%%\pacs[JEL Classification]{D8, H51}

%%\pacs[MSC Classification]{35A01, 65L10, 65L12, 65L20, 65L70}

\maketitle

%{ Hay que quitar el índice del preprint}
\tableofcontents

\section{Introduction}

The mathematical description of Quantum Field Theory on an arbitrary background has always been a complicated task. During the past century, most efforts focused on the algebraic side \cite{Haag1964, Streater1989,Doplicher1969,Doplicher1969a,Doplicher1971,Doplicher1974,Buchholz1982,Buchholz2020,Brunetti2015}. That approach proved to be very successful in encoding the properties of quantum operators and states in a mathematically consistent framework. Nonetheless, it did not consider other mathematical aspects of the description, as  the differential geometric ones; which were not described in detail either in the other main mathematical approach to QFT, non-commutative geometry (\cite{connesNoncommutativeGeometry1994,connesRenormalizationQuantumField2000}), more focused, again, on algebraic tools.   In particular, the geometrical description of the space of fields justifying, for instance, the existence of a Poisson bracket or  a symplectic form, relies on the definition of the suitable differentiable structure on the fields phase-space, which in this case is cumbersome to build. 

 Furthermore, one of our interests in this construction comes from its potential use in a hybrid quantum-classical Hamiltonian description of quantum fields dynamically coupled to the classical gravitational tensor fields \cite{alonso2023d}. The description of gravity largely depends on differential geometric tools. Thus, we tailor our  description of  the quantum scalar field  with those tools too.
 
 Regarding the symplectic characterization of quantum mechanics in  \cite{Kibble1979}, where Kibble defines a Poisson structure on the space of pure states of the Schrödinger-Hamiltonian description of  QFT. 
 In our work we consider the same idea and trace back the origin of different geometric tools  introduced on the phase spaces of fields and pure states. We aim to carefully study the interdependence of each approach to quantization with the complex structure of the classical phase space and, in turn, on the metric tensor defined on the space-time.  This aim justifies why we prefer  a non-covariant description of the quantum fields (although the global model is covariant) and the choice of the Hamiltonian 3+1 formalism for  gravity.      


Regarding the mathematical foundations, our construction is largely built upon the results of Hida-Malliavin calculus developed during the last two decades of the XXth century to study stochastic variational systems and applied with huge success to model financial systems \cite{Hida1980,Hida1993,Nunno2009, Wester2003}. These tools allow us to present the Hamiltonian descriptions of pure quantum states in a simplified language. The set of states has been usually studied with different algebraic or analytical tools under the names of  Schrödinger and holomorphic pictures \cite{Corichi2004,Oeckl2012,Oeckl2012a}.  In this two-part paper, we re-obtain those results but from a geometrical perspective which will allow us to build a description in terms of Poisson structures in the second part of this series  \cite{alonso2023a}.
  This geometrical flavor requires of a careful description of the different structures and of the relations between them required to build the final quantum description.  Hence, we want to build a consistent geometric description of  different sets of tools developed in the last years to model QFT. 

  In this first work we will consider the set of necessary mathematical tools to fix notation and provide the structures used in the second work of this series. The structure of the current paper is as follows. Firstly, with the objective of providing with the necessary framework to perform integro-differential calculus and differential geometry, in  section \ref{sec:math_preliminaries} we will briefly summarize how to use the structure of nuclear spaces  to build a differentiable structure on the set of fields. Following this discussion, the second mathematical ingredient within section \ref{sec:math_preliminaries} will delve on the necessity of rigged Hilbert spaces to identify these nuclear spaces with suitable subspaces of their duals to be able to map vectors and covectors on those field manifolds. This mechanism provides a practical framework to deal with tensor calculus.  A very important ingredient  that justifies the use of nuclear spaces to model the space will be  integral calculus. In particular the origin, existence and properties of the set of Gaussian measures on the space of fields, which constitute the third mathematical ingredient treated in section \ref{sec:math_preliminaries}. These Gaussian measures are required to build the Hilbert space structure on the set of quantum pure states.  In order to provide the reader with some insight on the nature of this Gaussian-measure based $L^2$ space, we present several functional dense subspaces of it of special relevance in QFT. 
 
  After this mathematical contextualization, we will be ready  to relate our framework with the usual approach to QFT, through the introduction of several mathematical transforms,  from stochastic calculus and Gaussian analysis. Firstly, the Wiener-Ito decomposition will yield the coefficients in the (bosonic) Fock space of quantum states. Secondly, the Segal-Bargmann transform will relate the holomorphic and Schrödinger functional pictures of QFT. Lastly,  we will introduce the main tools of Malliavin calculus, in particular the  role of  annihilation and creation operators in QFT will be played, respectively, by the  Malliavin derivative and the Skorokhod integral (generalization of Ito integral to non-adapted stochastic processes \cite{Nunno2009}).


This compendium of mathematical ingredients summits in section \ref{sec:Hida}, where we analyze the set of Hida test functions and their role in the rigorous characterization of the quantum pure states in the Schrödinger (and holomorphic) \textit{wave functional picture},  as dubbed in the physical literature.




Within the community of stochastic calculus, the relation between the tools of Malliavin calculus and operators in quantum field theory is already noticed in the literature \cite{henry-labordere2009} and  the close relation of white noise analysis and Hida-Malliavin calculus with Feynman integrals has been also studied  in detail \cite{Hida1993,Wester2003}.  From the physics point of view, the necessity of gaussian analysis has been also hinted as a necessary tool to study the Schrödinger picture of QFT \cite{Corichi2004}, the holomorphic picture \cite{Oeckl2012a} and their relation thorough the integral transforms \cite{Oeckl2012}. Moreover, gaussian analysis is needed to describe star products in  deformation quantization \cite{dito1990} and in algebraic QFT  quasifree states, also called gaussian states,  are considered the class of states that reproduce the desiderable physical properties of a theory \cite{Brunetti2015} and they can be easily related with the wick order of the Wiener-Ito decomposition theorem.

Despite the plethora of examples presented above that hits the importance of gaussian analysis as a tool for QFT, to our knowledge, the powerful and efficient formalism developed for the needs of stochastic calculus has not been used in full to describe the hamiltonian pictures of quantum field theory. Moreover, this formalism allows to take this analysis one step further and study, in a systematic way, the dynamics of the theory. The aim of this series of papers is to bring together both approaches to extract a coherent and systematic way  of treating QFT in curved spacetimes from the Hamiltonian perspective. 


Once the mathematical tools for the characterization of functional Hilbert spaces are introduced, section  \ref{sec:CFTQuantizar} analyzes the geometrical structure of the classical phase space of fields. As a result, we get that the introduction of a complex structure in that set  describes the Hilbert space of one-particle quantum states. We analyze also the role of the differential geometric properties of the phase space of classical fields with the construction of the quantum theory. This issue will be crucial when constructing the hybrid quantum-classical model. 

In section \ref{sec:conclusions} we summarize our results and contextualize the usefulness of this tools for applications in QFT. Precisely, in the second paper of this series \cite{alonso2023a}, we will use these tools to present different methods of quantization and characterize, in geometrical terms, the differentiable structures of the resulting manifold of quantum states. With these, we will obtain a geometrical description of quantum field theory on a curved background as a Hamiltonian system that will encode the definition of the dynamics from first principles.

\section{Mathematical preliminaries}
\label{sec:math_preliminaries}
In the following sections we will be constructing a differential-geometrical description of the phase spaces of classical and quantum field theories, which are, in both cases, infinite dimensional linear spaces. In order to do that, we will require modelling those manifolds on topological vector spaces where we can construct a differential calculus. Besides, in the case of the quantum phase space, we will also require the operators representing our physical magnitudes to be essentially self-adjoint on their domains, which will be submanifolds of the total space.  In some occasions, we may use a Hilbert space as modelling space, but this is not the general case. The most general framework is provided by dense subsets of functions, with enough properties to admit a well-behaved notion of continuity and differentiability. The best suited sets of such functions will be Nuclear-Fréchet (NF) spaces, whose main properties are presented below, with a few relevant examples. Furthermore, we will require to be able to define rigorously Borel probability measures on those manifolds, for the definition of quantum states to make sense so we are forced to consider also strong duals of NF spaces (DNF) as models of geometry. As gaussian measures are the most commonly used in QFT, this will be our preferred choice. 

Furthermore, the construction of tensors on these manifold will require considering both tangent and cotangent spaces. As the dual spaces of the nuclear spaces are spaces of distributions, tangent and cotangent spaces will be  non-isomorphic. In order to define a practical model, we will consider the notion of Gel'fand triple to be able to model vectors and covectors as elements of a suitable Hilbert space, which will provide us with suitable tensors on the tangent or on suitable dense subsets of the cotangent space. 

Finally, we will analyze in detail the nuclear space used to model the set of quantum pure states, the Hida test functions, and the corresponding Gel'fand triple. While, in regard to this paper, it is just an illustrative example of the construction for a certain nuclear space, its importance for the analysis in the second part of this series requires a specific analysis. 

%\subsection{Nuclear spaces and Gel'fand triples}

\subsection{Nuclear spaces: doing differential geometry in infinite dimensions}
It is a well known fact that, in order to define gaussian Borel probability measures in infinite dimensions, the natural setting is that of Fréchet Nuclear spaces (FN) and its strong duals (DFN). This is discussed thoroughly by \cite{Gellfand}.  Besides both categories of topological vector spaces are convenient and admit partitions of unity which turns them into suitable models for infinite dimensional geometry in the sense of \cite{Kriegl1997}.  In this section we focus on integration theory and we will summarize its main ingredients on the dual pair  $(\mathcal{N},\mathcal{N}')$ where $\mathcal{N}$ is a FN space and $\mathcal{N}'$ is its dual endowed with the strong topology\footnote{In this case we endow the DFN space with the strong topology, which coincides with the Mackey topology, but the same Borel $\sigma$-algebra is generated from the weak topology. }. In particular, we are interested in
$\big(C^{\infty}_c(K),D'(K)\big)$, the compactly supported functions and distributions over a compact subset $K\subset \mathbb{R}^{d}$ and $\big(\mathcal{S}(\mathbb{R}^d),\mathcal{S}'(\mathbb{R}^d)\big)$ the space of Schwartz functions of fast decay and tempered distributions. 
 
For completeness we define here the functional spaces to be used in a way that fits our purpose.

\begin{definition}
    \label{def:Nuclear}
    {\bf (Nuclear Space)} A nuclear space $X$ is a  topological vector space whose topology is given by a family of Hilbert seminorms  such that, if $X_{p}$ is the Hilbert space obtained as  the completion of the preHilbert  space $(X/p^{-1}(0), p)$, then there is a larger seminorm $q$ such that $X_q\hookrightarrow X_p$ is a Hilbert-Schmidt operator. 
\end{definition}

For infinite dimensional topological vector spaces $X,Y$ tensor products  can be endowed with a family of  different topologies between the projective $X\otimes_\pi Y$ and injective $X\otimes_\epsilon Y$ tensor product induced topologies.  If one or both spaces are nuclear both topologies coincide and therefore there is only one notion of tensor product $X\otimes Y$. We will  use this fact in following sections to define polynomials over $X$.

\begin{definition}
    \label{def:FN}
    {\bf (FN Fréchet Nuclear Space)} A Nuclear space which is also Fréchet  $\mathcal{N}$  can be regarded as the projective limit of a countable family of separable Hilbert spaces $\{\mathcal{H}_p\}_{p=0}^\infty$, with scalar products $\langle \cdot, \cdot  \rangle_p$ \ such that for $q>p$     
    the injection $\mathcal{H}_q\hookrightarrow \mathcal{H}_p$ is Hilbert-Schmidt. Thus
    $$
    \mathcal{N}=\bigcap_{p=0}^\infty\mathcal{H}_p 
    $$
    with the projective limit topology. 
\end{definition}

 Riesz representation theorem allows us to identify a Hilbert space with its dual. By convention we set $\mathcal{H}_0'=\mathcal{H}_0^*$ where $*$ represents complex conjugation. This implies that in the inclusion $\mathcal{H}_p\hookrightarrow\mathcal{H}_0$, for $p>0$, the space $\mathcal{H}_p$ can no longer be identified with its dual. We will denote this dual space by $\mathcal{H}_0\hookrightarrow\mathcal{H}_{-p}$. The next definition can be regarded as a consequence of the last one. 

\begin{definition}
    \label{def:DFN}
    {\bf (DFN Dual of a Fréchet Nuclear Space)} The strong dual of a FN space  $\mathcal{N}'$  can be regarded as the inductive limit of a countable family of Hilbert spaces $\{\mathcal{H}_{-p}\}_{p=0}^\infty$ such that for $q>p$
    the injection $\mathcal{H}_{-q}\hookleftarrow \mathcal{H}_{-p}$ is the transpose of a  Hilbert-Schmidt operator. Thus
    $$
    \mathcal{N}'=\bigcup_{p=0}^\infty\mathcal{H}_{-p} 
    $$
    with the inductive limit topology. A DFN space is always nuclear. Furthermore it is a Nuclear Silva space \cite{Kriegl1997}.
\end{definition}




\subsection{Rigged Hilbert spaces and main examples}

With the definitions given so far we notice that the space $\mathcal{H}_0$ plays an important role. This motivates the consideration of a triple of spaces therefore leading to the concept of rigged Hilbert space. 

\begin{definition} {\bf (Rigged Hilbert Space)}
    \label{def:RHS}
    Assume that $\mathcal{H}_0$ is a separable Hilbert space then, then we define a rigged Hilbert space or Gel'fand triple as the triple of spaces
\begin{equation*}
\label{eq:rigged}
\mathcal{N}\subset \mathcal{H}_0 \subset \mathcal{N}'
\end{equation*}
where $\mathcal{N}$ is a FN space and $\mathcal{N}'$ its dual and each inclusion is continuous. 

\end{definition}

In the structure here presented $\mathcal{H}_0$ is said to realize the duality, which means that the scalar product of $\mathcal{H}_0$ identifies an element of $\mathcal{N}$ with its dual in $\mathcal{N}'$. 
$\mathcal{N}$ is a dense subspace of the Hilbert space.
This must be seen as a way of parametrizing  $\mathcal{N}$ as a subset of $\mathcal{N}'$, in general different Hilbert spaces can be used to parametrize the duality of the same pair $(\mathcal{N},\mathcal{N}')$.



In the rigging there are three main ingredients. Firstly,  elements  $\xi_{\mathbf{x}}\in\mathcal{N}$ will be typically well defined functions (see examples below) and will be marked by a bold subindex.  Secondly,  elements $\varphi^{\mathbf{x}}\in\mathcal{N}'$ will be distributions with no pointwise meaning in general and they will be marked by a bold superindex. Sometimes we will denote the value of the function at a point $x$  with  a regular subindex $\xi_x$, this can't be done in general for distributions. With those two ingredients the duality is denoted, mimicking Einstein summation convention, 
\begin{equation}
\label{eq:dualityNotationDistributions}
\xi_{\mathbf{x}}\in\mathcal{N}, \varphi^{\mathbf{x}}\in\mathcal{N}' \textrm{ paired by } \xi_x\varphi^x :=\langle\varphi^{\mathbf{x}},\xi_{\mathbf{x}}\rangle \in \mathbb{R}\ \ (\mathbb{C}).
\end{equation}

The third and last ingredient is the dual pairing of $\mathcal{H}_0$. This element can be seen in a threefold way
\begin{enumerate}[label=\alph*)]
    \item Dual pairing between elements in $\mathcal{N}\subset \mathcal{H}_0$ with $\langle\xi^1_{\mathbf{x}} ,\xi^2_{\mathbf{y}}\rangle =\xi^1_x\Delta^{xy}\xi_y^2$
    \item Dual pairing between elements in  $ \mathcal{H}_0\subset \mathcal{N}'$ with $\langle\varphi_1^{\mathbf{x}} ,\varphi_2^{\mathbf{y}}\rangle =\varphi_1^xK_{xy}\varphi^y_2$ 
    \item Injection of ${\Delta^{\mathbf{xy}}}:\mathcal{N}\hookrightarrow \mathcal{N}'$ trough the pairing of $\mathcal{H}_0$ such that $\xi_x\Delta^{xy}=\xi^y$ with a densely defined inverse $\varphi^xK_{xy}=\varphi_y$ such that $\Delta^{xz}K_{zy}=\delta^x_y$ is the evaluation mapping.
\end{enumerate}


The main examples of FN spaces, which play the roles of model spaces for the classical field theories in this work are the following.

\begin{example}\label{ex:testFunctions}{\bf (Smooth functions of compact support on a compact $K$)}
    Let $K\subset \mathbb{R}^d$ then the space $C^{\infty}_c(K)$ is FN with the initial topology induced from the family of Hilbertian norms

\begin{equation}
\label{eq:hilbertianNormsCompactSupport}
\lVert\varphi\rVert_{K,n}^2=\sum_{\lvert\alpha\rvert\leq n}\int_{K}\lvert\partial^\alpha\varphi(x)\rvert^2 dx^d,
\end{equation}
where $\alpha\in \mathbb{N}_0^d$ is a multi-index and $\lvert\alpha\rvert=\alpha_1+\cdots+\alpha_d$.  Its dual is the DFN space $D'(K)$. 

Notice that we could also define $C^\infty_c(U)$, for $U\subset \mathbb{R}$ an open set. This is a nuclear  but not  a Fréchet space; instead it is a Nuclear LF space (see \cite{Gellfand} for details). For our purposes we will limit our attention to the compact domain in order to reduce  technicalities.
\end{example}



\begin{example}{\bf (Smooth functions of fast decay)}\label{ex:schwartFunctions} The space $\mathcal{S}(\mathbb{R}^d)$ of smooth functions of fast decay is FN with the initial topology induced from the family of Hilbertian norms

    \begin{equation}
        \label{eq:hilbertianNormsSchwartz}
        \lVert\varphi\rVert_{m,n}^2=\sum_{\lvert\alpha\rvert\leq n, \lvert\beta\rvert\leq m}\int_{\mathbb{R}^d}\prod_{i=1}^d(1+x_i^2)^{2\beta_i}\lvert \partial^\alpha\varphi(x)\rvert^2 dx^d
    \end{equation}
    where $\alpha,\beta\in \mathbb{N}_0^d$ are multi-indices. This topology is equivalent to the one usually introduced via supremum type seminorms. Its dual $\mathcal{S}'(\mathbb{R}^d)$ is the DFN space of tempered distributions. This space is the one commonly used in white noise analysis \cite{Hida1993}  
\end{example}


\begin{example}\label{ex:testFunctionsManifold}{\bf (Smooth functions of compact support on a compact Manifold )} For $\Sigma$ a compact manifold the definition of $C^\infty_c(\Sigma)$ goes on as in \autoref{ex:testFunctions}. If we endow the manifold with a riemannian metric $h$ there is a canonical volume form 
\begin{equation}
\label{eq:riemannianVolumeform}
\int_\Sigma f\mathrm{dVol}_h=\int_U dx^d \sqrt{\lvert h(x)\rvert} f(x),
\end{equation}
for $f$ supported on a subset of $\Sigma$ covered by  a single chart with coordinates in $U$. This, in turn, leads to a canonical Hilbert space of square integrable functions $L^2(\Sigma,\mathrm{dVol}_h)$ and therefore to the rigged Hilbert space 
\begin{equation}
\label{eq:riggedOnAManifold}
C^\infty_c(\Sigma)\subset L^2(\Sigma,\mathrm{dVol}_h) \subset D'(\Sigma)
\end{equation}  
  

As discussed in \autoref{def:RHS} the scalar product on a rigged Hilbert space can be considered in a threefold way, as we can illustrate over a chart with coordinates in $U\subset \mathbb{R}^d$

\begin{enumerate}[label=\alph*)]
        \item Dual pairing between elements in $C^\infty_c(\Sigma)\subset L^2(\Sigma,\mathrm{dVol}_h)$ with $$\langle f_\mb{x} ,g_\mb{y}\rangle =f_x\delta^{xy}g_y=\int_U dx^d\sqrt{\lvert h(x)\rvert} f(x)g(x)$$
    \item Dual pairing between elements in  $L^2(\Sigma,\mathrm{dVol}_h) \subset D'(\Sigma)$ with $$\langle\phi^{\mathbf{x}} ,\varphi^{\mathbf{y}}\rangle =\phi^x\delta_{xy}\varphi^y=\int_U\frac{dx^d}{\sqrt{\lvert h(x)\rvert}} \phi(x)\varphi(x),$$
    where $\phi(x),\varphi(x)$ are identified with  densities of weight 1. Notice that the factor $\sqrt{\lvert h(x)\rvert}$ appears in the denominator and $\frac{\phi(x)}{\sqrt{\lvert h(x)\rvert}}, \frac{\varphi(x)}{\sqrt{\lvert h(x)\rvert}}$ are $L^2(\Sigma, \mathrm{dVol}_h)$ representatives.
    \item We can also see the injection ${\delta^{\mathbf{xy}}}:C^\infty_c(\Sigma) \hookrightarrow D'(\Sigma)$  such that
    $$f^y=f_x\delta^{xy}=\int_U dx^d\sqrt{\lvert h(x)\rvert}\delta^d(x-y) f(x)=\sqrt{\lvert h(y)\rvert}f(y)$$
    is a representative of a density of weight 1. On the other hand $\delta_{\mathbf{xy}}$ is a densely defined inverse 
    $$\varphi^y=\varphi^x\delta_{xy}=\int_U \frac{dx^d}{\sqrt{\lvert h(x)\rvert}}\delta^d(x-y) \varphi(x)=\frac{\varphi(y)}{\sqrt{\lvert h(y)\rvert}}$$
    such that $\delta^{xz}\delta_{zy}=\delta^x_y$ is the evaluation mapping
    $$\delta^x_yf_x=f_y=f(y)$$
    which only makes sense over $C^\infty_c(\Sigma)$.
\end{enumerate}

The differences between $\delta^{\mathbf{xy}}$, $\delta_{\mathbf{xy}}$ and $\delta^{\mathbf{x}}_{\mathbf{y}}$ are often misinterpreted in the literature since the three of them can be identified with the Dirac delta in $\mathbb{R}^d$ with the Lebesgue measure. 

In this picture elements of  $D'(\Sigma)$ are distributions used to integrate the elements of $C^\infty_c(\Sigma)$ and this nontrivial structure is reframed in the rigged Hilbert space with a nontrivial representation of the scalar product as an integral kernel.

For different applications, the identification of the field and its dual may be different. Thus, in  classical field theory the common approach \cite{gotayMomentumMapsClassical2004a,gotayMomentumMapsClassical2004} is to model classical fields as elements of a $n$-th jet bundle with $n$ the number of derivatives needed to describe the theory. We can say that, in this approach, fields are modelled over elements of  $C^\infty_c(\Sigma)$. In those approaches the dual jet bundle is built upon the concept of density of weight one  that are constructed using volume forms $f dx^1\wedge \cdots\wedge dx^d$ with $f\in C^\infty_c(\Sigma)$.  Thus, the topology of the model space for the dual jet bundle is that of  $C^\infty_c(\Sigma)$. This approach makes the treatment significantly easier at the cost of giving up the possibility of describing the duality with the strong topology. In those approaches the dual pairing is described as in a).

In algebraic QFT the situation is radically different, classical fields to be quantized are taken as a priori distributions of $D'(\Sigma)$ then both fields and their duals are modelled over the same space and the dual pairing only makes sense as in case b) \cite{Brunetti2015}. 

\end{example}


Other important family of examples of rigged Hilbert spaces are reliant upon the FN space of  Hida test functions, which will be the set used to model the quantum field pure state. Its construction is significantly more difficult and needs a careful study of gaussian integration to be developed in the following sections. For this reason  we postpone its introduction to \autoref{sec:Hida}. 

\subsection{Defining a scalar product: Gaussian measures}
As we are intending to create a quantum description of our physical state, we must endow the space of quantum states with a scalar product.  As we mentioned above, the most commonly used measures in QFT correspond to gaussian measures, often referred to a given vacuum state. 
We will review below their main definition and  properties, first for the simple case of real test functions, and then in the case of complex ones, which will be more useful for us later.
\subsubsection{The real case}
The procedure to define a Gaussian measure over $\mathcal{N}'$ is straightforward once we notice that the $\sigma$-algebra generated by its cylinder subsets corresponds to the Borel $\sigma$-algebra of the weak and strong topologies of $\mathcal{N}'$.  The only ingredient needed to define such measure is the characteristic functional,
\begin{equation}
    \label{eq:chartacteristicFunctional}
        C(\xi_{\mathbf{x}})=\int_{\mathcal{N}'}D\mu(\varphi^{\mathbf{x}}) e^{i\xi_x\varphi^x} =e^{-\frac{1}{2}{\xi}_x\Delta^{xy}\xi_y},
\end{equation}
 Minlos theorem \cite{Gellfand,Hida1993} ensures that, if ${\xi}_x\Delta^{xy}\xi_y$ is continuous with respect to the topology of $\mathcal{N}$, there is a unique Borel probability measure that fulfills \eqref{eq:chartacteristicFunctional}. Considered with respect to a Gel'fand triple,
$$
\mathcal{N}\subset \mathcal{H}_\Delta \subset \mathcal{N}',
$$
the bilinear symmetric form $\Delta$ corresponds to the representation of the scalar product of the Cameron-Martin Hilbert space $\mathcal{H}_\Delta$. This space is obtained by completion of $\mathcal{N}$ with the covariance $\Delta^{\mathbf{xy}}$ and is of huge importance in gaussian analysis as we will see below. 


\subsubsection{The complex case}
If we consider the case of complex test functions $\mathcal{N}_\mathbb{C}$ of any of the aforementioned examples of nuclear spaces, most of the properties work in the same form as for the real case.  Let us consider the complexification of the space of distributions such that  an element $\phi^{\mathbf{x}}\in\mathcal{N}'_\mathbb{C}$ is written as $\phi^{\mathbf{x}}=\frac{1}{\sqrt{2}}(\varphi^{\mathbf{x}}+i\pi^{\mathbf{x}})$ for $(\varphi^{\mathbf{x}},\pi^\mathbf{x})\in \mathcal{N}'\times \mathcal{N}'$. Complex test functions can also be written in holomorphic and antiholomorphic  coordinates, for $\rho_{\mathbf{x}}\in\mathcal{N}_\mathbb{C}$ it follows  $\rho_x=\frac{1}{\sqrt{2}}(\xi^{\varphi}_x+i\xi^{\pi}_x)$ and $\bar\rho_x=\frac{1}{\sqrt{2}}(\xi^{\varphi}_x-i\xi^{\pi}_x)$. If we endow each copy of the cartesian product with the measure given by \eqref{eq:chartacteristicFunctional},  we can define a measure to the complex domain  defined by the characteristic functional  

\begin{equation}
\label{eq:complexCharacteristicFunctional}
C_c(\rho_{\mathbf{x}},\bar\rho_{\mathbf{y}})=\int_{\mathcal{N}_{\mathbb{C}}'}D\mu_c(\bar \phi^{\mathbf{x}}, \phi^{\mathbf{x}}) e^{i({\rho_x\bar{\phi}^x+\bar\rho_x{\phi}^x})} =e^{-\bar{\rho}_x\Delta^{xy}\rho_y},
\end{equation}
where now the covariance $\Delta^{xy}$ will be a hermitian form on $\mathcal{N}_\mathbb{C}$, i.e. $\Delta^{xy}=(\Delta^*)^{yx}$. 
Notice that the complex measure is the product measure of two real ones.

\subsubsection{Integral transforms on gaussian measures}
Among the different integral transforms defined over a gaussian measure the most relevant for Quantum Field Theory is the $S$ transform. Throughout this whole section we will consider a gaussian measure $D\mu$ defined on $\mathcal{N}'$ or $\mathcal{N}'_\mathbb{C}$ with associated characteristic functional \eqref{eq:complexCharacteristicFunctional}. For the real case we define:

\begin{definition}
    \label{def:Stransform}{\bf ($S_\mu$ transform)} For the  gaussian measure of \eqref{eq:chartacteristicFunctional} we define the $S_\mu$ transform as a generalization of the Laplace transform. For a given $\Psi:\mathcal{N}'\to \mathbb{C}$ it takes the form 
    \begin{equation}
        \label{eq:stransform}
        S_\mu[\Psi](\xi_\mathbf{x})=\int_{\mathcal{N}'}D\mu(\varphi^{\mathbf{x}})\Psi(\varphi^{\mathbf{x}})\exp{\big(\xi_x\varphi^x-\frac12\xi_x\Delta^{xy}\xi_y\big)}
\end{equation}
\end{definition}
This is to be thought as the action of a translation $\varphi^{\mathbf{x}}\to\varphi^{\mathbf{x}}+\Delta^{\mathbf{x}y}\xi_y$ on the domain of integration . For the complex case it is readily defined in (anti)-holomorphic coordinates as 

\begin{definition}
    \label{def:SComplextransform}{\bf ($S_{\mu_c}$ complex transform)} For the complex gaussian measure of \eqref{eq:complexCharacteristicFunctional} and a given $\Psi:\mathcal{N}_{\mathbb{C}}'\to \mathbb{C}$, the transform takes the form 
    \begin{equation}
        \label{eq:stransformcomplex}
        S_{\mu_c}[\Psi](\rho_\mathbf{x},\bar\rho_\mathbf{x})=\int_{\mathcal{N}'_{\mathbb{C}}}D\mu_c(\phi^{\mathbf{x}})\Psi(\phi^{\mathbf{x}},\bar\phi^{\mathbf{x}})\exp{\big({\rho}_x\bar\phi^x+\bar{\rho}_x{\phi}^x-\bar\rho_x\Delta^{xy}\rho_y\big)}
        \end{equation}
\end{definition}

Notice that this definition is nothing but the combination of two copies of \autoref{def:Stransform}, one for each variable, and the required change to holomorphic coordinates.

Another relevant transform is the $T$ transform. It is the  extension of  the definition of characteristic functional \eqref{eq:complexCharacteristicFunctional} to  include a generic function besides the measure, defined as follows
\begin{definition}
    \label{def:TComplextransform}{\bf ($T_{\mu_c}$ complex transform)} For the complex gaussian measure of \eqref{eq:complexCharacteristicFunctional} and a given $\Psi:\mathcal{N}_{\mathbb{C}}'\to \mathbb{C}$  takes the form 
\begin{equation}
    \label{eq:Ttransform}
    T_{\mu_c}[\Psi](\rho_\mathbf{x},\bar{\rho}_\mathbf{y})=\int_{\mathcal{N}'_\mathbb{C}}D\mu_c(\phi^{\mathbf{x}})e^{i( \rho_z\bar\phi^z+\bar\rho_z\phi^z)} \Psi(\phi^{\mathbf{x}},\bar\phi^\mathbf{y})
    \end{equation}
\end{definition}

\noindent Both transforms are related by $C(\rho_\mathbf{x},\bar{\rho}_\mathbf{y}) S_{\mu_c}[\Psi](\rho_\mathbf{x},\bar\rho_\mathbf{x})= T_{\mu_c}[\Psi](-i\rho_\mathbf{x},-i\overline{\rho}_\mathbf{y})$ each of those transformations reveal  different features of $L^2(\mathcal{N}'_\mathbb{C},D\mu_c)$ as we will see bellow.  


\subsection{Dense subsets of $L^2(\mathcal{N}'_\mathbb{C},D\mu_c)$}

The structure of the  Hilbert space $L^2(\mathcal{N}'_\mathbb{C},D\mu_c)$ described in the previous subsection may seem difficult to tackle down. At first glace we are dealing with functions with domains in an infinite dimensional space of distributions. This is some times referred to as working at the level of {\it second quantization} in the literature. In order to study the whole space, it is useful to consider some simpler to define families of functions that are dense in the space. Here we consider the set of polynomials, trigonometric exponentials and coherent sates. We claim that they are dense subsets of the whole space, for a proof of this statement in each case we refer to- general arguments made at the beginning of  \cite{Hida1993}. Each of those families are easily and unambiguously defined thanks to the properties of nuclear spaces and each of them show different properties when combined with the S and T transforms introduced in the previous section. 

\subsubsection{The algebra of Polynomials}
\label{sssec:Polynomials}

 The most important dense subset in gaussian analysis are polynomials.  The fact tha $\mathcal{N}'_{\mathbb{C}}$ is a nuclear space implies that the tensor product $ {\mathcal{N}_\mathbb{C}'}^{\otimes k}$ is unambiguously defined as well as $\mathcal{N}_{\mathbb{C}}^{\otimes n}$ for $k,n$ positive integers. With this remark we define the algebra of  polynomials $\mathcal{P}(\mathcal{N}'_\mathbb{C})$, with coefficients in $\mathcal{N}_{\mathbb{C}}$, as the set of polynomials of arbitrary degrees in which  a polynomial of degrees $(n,m)$ is  written as


 \begin{equation}
     \label{eq:polynomialscompl} 
     P_{nm}(\phi^{\mathbf{x}},\bar\phi^{\mathbf{x}})= \sum_{k=0}^n\sum_{l=0}^m C^{kl}_{\vec{x}_k,\vec{y}_l}(\phi^k)^{\vec{x}_k}(\bar\phi^l)^{\vec{y}_l}
 \end{equation}
 With $C^{kl}_{\vec{\mathbf{x}}_k,\vec{\mathbf{y}}_l}\in \mathcal{N}_{\mathbb{C}}^{\otimes k+l}$ and $(\phi^k)^{\vec{\mathbf{x}}_k}\in {\mathcal{N}_\mathbb{C}'}^{\otimes k}$, $(\bar\phi^l)^{\vec{\mathbf{x}}_l}\in {\mathcal{{N}_\mathbb{C}'}^*}^{\otimes l}$ are the tensor product of copies of the same distribution. Here we introduce the notation $\varphi_{\vec{x}_n}:=\varphi_{x_1,\cdots,x_n}$.
 {\color{blue} To make sense of this expression is enough to consider a coefficient $C^{kl}_{\vec{\mathbf{x}}_k,\vec{\mathbf{y}}_l}=\bar\xi^1_{x_1}\cdots\bar\xi^k_{x_k}{\eta}^1_{y_1}\cdots{\eta}^l_{y_l}$ because these coefficients generate the whole space. Then, using the natural dual pairing of $\mathcal{N}_{\mathbb{C}}$ with its dual expressed as $\langle \alpha,\phi \rangle= \bar\alpha_{x}\phi^x$, we get 
 $$
 C^{kl}_{\vec{x}_k,\vec{y}_l}(\phi^k)^{\vec{x}_k}
 (\bar\phi^l)^{\vec{y}_l}=\prod_{i=1}^k 
 \langle \xi_i,\phi \rangle
 \prod_{j=1}^l
  \overline{\langle \phi, \eta_j \rangle}
$$
 }
 Thus polynomials defined in \eqref{eq:polynomialscompl}  have pointwise meaning because each dual pairing is finite. Nonetheless, there is an ambiguity remaining in this definition  because the monomials are contractions of  $n$ or $m$ products of the same distribution. This fact implies there are different coefficients that provide the same polynomial. To kill this ambiguity we introduce the symmetrization operator 
    \begin{equation}
    \label{eq:symmetrization}
    \varphi_{(x_1,\cdots,x_n)}:=\frac{1}{n!}\sum_{\pi\in S_n}\varphi_{x_{\pi(1)},\cdots,x_{\pi(n)}}
    \end{equation}

 Then for polynomials to be unambiguously defined through its coefficients we ask for them to fulfill $C^{kl}_{(\vec{\mathbf{x}}_k),(\vec{\mathbf{y}}_l)}=C^{kl}_{\vec{\mathbf{x}}_k,\vec{\mathbf{y}}_l}$ and we call them bosonic coefficients. Among polynomials we can find an orthonormal basis of the whole space and then describe it efficiently through its bosonic coefficients in an appropriate way using the $S$ transform. This feature is so rich and enlightening that it deserves its own separate analysis that will be the subject of study of \autoref{ssec:WienerIto}. 
 
\subsubsection{ The algebra of trigonomeric exponentials: The nature of Reproducing Kernel Hilbert space}

 It is sometimes useful to consider functions outside polynomials. For that reason we turn our attention to the set of trigonometric exponentials $\mathcal{T}(\mathcal{N}'_\mathbb{C})$, which can be seen  to be a dense subset of $L^2(\mathcal{N}_\mathbb{C}, D\mu_c)$. They are defined as   
\begin{equation}
\label{eq:trigExponentials}
\mathcal{T}(\mathcal{N}'_\mathbb{C})=\left\{\textrm{ Finite linear combinations of }  \mathcal{E}_{\chi}=\exp\left(i\bar\chi_x\phi^x+i\chi_x\bar\phi^x\right)\right\}
\end{equation}
where $\chi_{\mathbf{x}}\in \mathcal{N}_{\mathbb{C}}$. Below we perform an analysis from the point of view of a basis of trigonometric exponentials. The tool at hand  to carry on with our survey is the $T_{\mu_c}$ transform of \autoref{def:TComplextransform}, its action over $\mathcal{T}(\mathcal{N}'_\mathbb{C})$ is completely obtained by linearity from  
\begin{equation}
\label{eq:TtransformTrig}
T_{\mu_c}[\mathcal{E}_{\chi}](\rho_\mathbf{x},\bar{\rho}_\mathbf{y})=C_c(\rho_{\mathbf{x}}+\chi_{\mathbf{x}},\overline{\rho_{\mathbf{y}}+\chi_{\mathbf{y}}}).
\end{equation}
Here $C_c$ is the characteristic functional given by \eqref{eq:complexCharacteristicFunctional}. Notice that in this equation $\rho_\mathbf{x}$ plays the role of a variable while $\chi_\mathbf{x}$ is a parameter. Hence, we interpret this expression as a function of $\rho_\mathbf{x}$ and extend the domain of  $T_{\mu_c}[\mathcal{E}_{\chi}](\rho_\mathbf{x},\bar{\rho}_\mathbf{y})$ to make sense in the whole Cameron-Martin Hilbert space  $\rho_{\mathbf{x}}\in \mathcal{H}_\Delta$. The next step is to construct a reproducing kernel Hilbert space $\mathcal{R}$. In order to do so, we firstly introduce the kernel functional defined as

\begin{equation}
    \label{eq:kernelforRKHS}
    \begin{array}{cccc}
    RK(\cdot,\cdot) :&\mathcal{H}_\Delta \times \mathcal{H}_\Delta & \longrightarrow & \mathbb{R} \\ 
     & (\rho_\mathbf{x},\chi_\mathbf{y}) &\longmapsto & C_c(\chi_{\mathbf{x}}-\rho_{\mathbf{x}},\overline{\chi_{\mathbf{y}}-\rho_{\mathbf{y}}}).
    \end{array}
    \end{equation}

    Using this definition, $\mathcal{R}$ is defined from the space of linear combinations of  functions (of $\rho,\bar\rho$) spanned by the r.h.s of \eqref{eq:TtransformTrig} as its completion with a sesquilinear scalar product $\langle,\rangle_{\mathcal{R}}$. It is enough to define this scalar product through its action over trigonometric exponentials under $T_{\mu_c}$ as

\begin{equation}
\label{eq:scalarProductCoherent}
\langle T_{\mu_c}[\mathcal{E}_{\rho}] , T_{\mu_c}[\mathcal{E}_{\chi}]\rangle_\mathcal{R}=RK(\rho_{\mathbf{x}},\chi_{\mathbf{x}}).
\end{equation}

A key factor in this construction is that $RK(\rho_{\mathbf{x}},\bullet )\in \mathcal{R}$ and for every $F\in \mathcal{R}$ we have  $\langle RK(\rho_{\mathbf{x}}, \bullet  ) , F(\bullet)\rangle_\mathcal{R}= F(\rho_{\mathbf{x}})$, this is, point evaluation is continuous in $\mathcal{R}$ which is the definition of reproducing Kernel Hilbert space.  

The remaining part of the discussion is to show that $T_{\mu_c}:L^2\big({\mathcal{N}_{\mathbb{C}}'},D\mu_c\big) \to \mathcal{R}$ is an isometric isomorphism of Hilbert spaces. This easily follows from 

\begin{equation}
\label{eq:ISometry}
\int_{\mathcal{N}'_\mathbb{C}}D\mu_c\overline{ \mathcal{E}_{\chi}} \mathcal{E}_{\rho}
=\langle T_{\mu_c}[\mathcal{E}_{\chi}] , T_{\mu_c}[\mathcal{E}_{\rho}]\rangle_\mathcal{R}
\end{equation}


From this perspective we may ask in which sense point evaluation is allowed for elements in $L^2\big({\mathcal{N}_{\mathbb{C}}'},D\mu_c\big)$. It is important to notice that the domain of the functions that form this space is much larger than that of the functions of $\mathcal{R}$, given by $\mathcal{H}_\Delta$. Therefore there is no hope for the full notion of reproducing kernel  to be found in the former as we did in the latter, and as such, there is no notion of point evaluation to be found. The issue of point evaluation in this case is similar to the issue of defining a Dirac 
delta in finite dimensional domains. As in that case there is a notion of evaluation that makes sense in a distributional sense as we will see in \autoref{ex:repKernel} below.  At the level of Hilbert space, the family of functions considered in the next section provides tools to delve further in this question.


 
\subsubsection{The algebra of coherent states: Subspace of Holomorphic functions}

In this case we combine features from the two families considered above. From the algebra of polynomials we notice that there is a decomposition in holomorphic and antiholomorphic functions as follows

\begin{equation}
    \label{eq:ComplexHilbertspacedecomp}
    L^2\big({\mathcal{N}_{\mathbb{C}}'},D\mu_c\big)= L^2_{Hol}\big({\mathcal{N}_{\mathbb{C}}'}, D\mu_c\big)\bar{\otimes}\overline{ L^2_{{Hol}}\big({\mathcal{N}_{\mathbb{C}}'}, D\mu_c\big)}
    \end{equation}
    where $\bar{\otimes}$ represents the topological tensor product completed with the natural Hilbert topology induced by the tensor product. Plus, an element $\Psi_{Hol}\in L^2_{Hol}\big({\mathcal{N}_{\mathbb{C}}'},D\mu_c\big) $ is decomposed as    \begin{equation}
        \label{eq:chaosdecompositionhol}
        \Psi_{Hol}(\phi^{\mathbf{x}})=\sum_{n=0}^\infty\psi^{n}_{\vec{x}_n}(\phi^n)^{\vec{x}_n}.
    \end{equation}
    Here we allow the bosonic coefficients $\psi^{n}_{\vec{x}_n}$ to be elements of a larger space $\mathcal{H}_\Delta^{\otimes n}$. This is a general feature of the Wiener-Ito decomposition theorem that will be treated in detail in \autoref{ssec:WienerIto}.


Now we particularize our analysis to the subspace $L^2_{Hol}\big({\mathcal{N}_{\mathbb{C}}'},D\mu_c\big)$. Besides holomorphic polynomials, this space  contains another dense subspace of interest. We introduce the algebra of functions generated by (Holomorphic) coherent states $\mathcal{K}_{\chi}(\phi^{\mathbf{y}})=e^{\bar\xi_x\phi^x}$ and we denote it 
\begin{equation}
    \label{eq:CoherentStates}
    \mathcal{C}_{Hol}(\mathcal{N}'_\mathbb{C})=\left\{\textrm{ Finite linear combinations of }  \mathcal{K}_{\chi}=\exp\left(\bar\chi_x\phi^x\right)\right\}
    \end{equation}

Based on this subset calculus is easily conducted on the whole space by just a few rules.  First, notice that
the scalar product of two coherent states corresponds to
\begin{equation*}
    \langle \mathcal{K}_{\chi},\mathcal{K}_{\rho}\rangle = e^{{\bar{\chi}_x\Delta^{xy}\rho_y}},
\end{equation*}
from which, by linearity, stems the reproducing property
\begin{equation*}
    \langle \mathcal{K}_{\chi},\Psi\rangle= \Psi(\chi_{x}\Delta^{x \mathbf{y}}). \\
\end{equation*}
 Hence $ \mathcal{K}_{\chi}(\phi^{\mathbf{y}})\in L^2_{Hol}\big({\mathcal{N}_{\mathbb{C}}'},D\mu_c\big)$. Those coherent states can be extended to be parametrized by elements $\rho^{\mathbf{x}}\in \mathcal{H}_\Delta$ regarded as distributions with $ \mathcal{K}_\rho(\phi^{\mathbf{y}})=e^{\bar\rho^xK_{xy}\phi^x}$, the duality expressed as in option b) of the previous section in terms of the inverse of the covariance, $K_{\mathbf{xy}}$, and therefore, lacking any pointwise meaning. 
 
 
 This subspace possesses good analytical properties similar to those of holomorphic functions in finite dimensions. We can think about it as the Hardy space of square integrable holomorphic functions  in an infinite dimensional space whose orthogonal directions are provided by the covariance $\Delta$. To stablish this relation in detail let $Hol(\mathcal{H}_\Delta)$ be the space of holomorphic functions over the Hilbert space $\mathcal{H}_\Delta$ defined by

\begin{definition}
    \label{def:Holomorphic}{\bf (Holomorphic function)}
    Let $\Psi:\mathcal{H}_\Delta\to\mathbb{C}$ be a function defined on an open set $U$. $\Psi$ is said to be G-Holomorphic on $U$ if for all $\phi^\mathbf{x}_0,\phi^\mathbf{x}\in U$ the function
    \begin{equation}
    \label{eq:complexSubfunctionGateaux}
    \begin{array}{rccc}
        \Psi_{\phi^\mathbf{x}_0,\phi^\mathbf{x}}:&\mathbb{C}&\longrightarrow&\mathbb{C}\\
         & z & \longmapsto & \Psi(\phi^\mathbf{x}_0+z\phi^\mathbf{x})
    \end{array}
    \end{equation} 
    is holomorphic. 

    As usually happens in infinite dimensions, continuity does not follow from differentiability and we must ensure the continuity. Thus $\Psi$ is said to be Holomorphic on $U$ if it is G-Holomorphic and locally bounded.
    \end{definition}
    

We endow $Hol(\mathcal{H}_\Delta)$  with the topology of convergence on bounded subsets. It follows  
\begin{equation}
\label{eq:injectionHolomorphic}
\langle \mathcal{K}_{\rho},\cdot\rangle: L^2_{Hol}\big({\mathcal{N}_{\mathbb{C}}'},D\mu_c\big) \to H^2(\mathcal{H}_\Delta)\subset Hol(\mathcal{H}_\Delta) 
\end{equation}
 is an injective continuous linear mapping \cite{Oeckl2012a}. Where $H^2(\mathcal{H}_\Delta)$ stands for the image of this map. Moreover, the reproducing kernel 
\begin{equation}
\label{eq:kernel}
RK_{Hol}(\chi_{\mathbf{x}},\rho_\mathbf{y})=e^{\overline{\chi}_x\Delta^{xy}\rho_y}
\end{equation}
turns  $(H^2(\mathcal{H}_\Delta),RK_{Hol})$ into a reproducing kernel Hilbert space and \eqref{eq:injectionHolomorphic} stablish an isometric isomorphism.

A similar discussion can be carried out for the anti-holomorphic functions but in that case the isomorphism is realized by 

\begin{equation}
\label{eq:injectionAntiHolomorphic}
\langle \overline{\mathcal{K}_{\rho}},\cdot\rangle: L^2_{\overline{Hol}}\big({\mathcal{N}_{\mathbb{C}}'},D\mu_c\big) \to \overline{H^2}(\mathcal{H}_\Delta)
\end{equation}

Finally, the $T_{\mu_c}$ transform of \autoref{def:TComplextransform} is written  $T_{\mu_c}=\langle \overline{\mathcal{K}_{i\rho}}\otimes \mathcal{K}_{-i\rho} ,\cdot\rangle$. Therefore, there is an isometric isomorphism of reproducing kernel Hilbert spaces such that

\begin{equation}
\label{eq:isomorphism}
    H^2(\mathcal{H}_\Delta)\otimes \overline{H^2}(\mathcal{H}_\Delta) \cong \mathcal{R}
\end{equation}

Where $\mathcal{R}$ is the reproducing kernel Hilbert space described in the previous section.

\subsubsection{Brief physical implications}

So far we have seen how the dense subsets of trigonometric exponentials and coherent states unveil several properties of the space $L^2\big({\mathcal{N}_{\mathbb{C}}'},D\mu_c\big)$. In \cite{alonso2023a} this kind of  space will play a twofold role. On one hand a gaussian Hilbert space of complex domain is a prequantum Hilbert space that is a crucial step in the geometric quantization procedure. On the other a similar gaussian space (with different covariance), represents the class of classical functions that can be quantized under Weyl quantization.

Among these properties the most important one for geometric quantization is the decomposition in holomorphic and antiholomorphic parts \eqref{eq:ComplexHilbertspacedecomp}. This is because, as we will see in \cite{alonso2023a}, the Hilbert space of pure states in the quantum field theory can be represented in several equivalent ways from which the holomorphic representation selects the subspace  $L^2_{Hol}\big({\mathcal{N}_{\mathbb{C}}'},D\mu_c\big)$.
Finally, the fact that \eqref{eq:injectionHolomorphic} provides a isomorphism of Hilbert spaces with  $H^2(\mathcal{H}_\Delta)$, a reproducing kernel Hilbert space with a well behaved domain $\mathcal{H}_\Delta$, ensures nice analytical properties for  $L^2_{Hol}\big({\mathcal{N}_{\mathbb{C}}'},D\mu_c\big)$ that allows to perform efficient calculus and study quantum field theory keeping technicalities under control.

For the case of Weyl quantization in \cite{alonso2023a} we show that it is enough to understand the quantization of trigonometric exponentials \eqref{eq:trigExponentials}  as plane waves of the classical field theory and the machinery of reproducing kernels allows to quantize them according to the Weyl or canonical commutation relations. 

The features unveiled by the family of polynomials are much richer and we devote the next section to them. We advance that the Wiener-Ito decomposition theorem will be the tool that provides the particle interpretation of the QFT.


\subsection{Modelling quantum field pure states I: Wiener Ito decomposition}
\label{ssec:WienerIto}
In this section we will consider the mathematical tools required to represent the concept of pure state of a quantum field and its main properties, in particular the relation with the concept of physical particles.

For a scalar filed we can consider two types of unitary equivalent representations.  On one hand, the Hilbert space of the Schrödinger wave functional representation $\Psi[\varphi]$ of quantum states is given by $L^2(\mathcal{N}', D\mu)$, for a given Gaussian measure $D\mu(\varphi)$, which, in the physics literature, is associated to the vacuum state of the theory. On the other hand, for the Holomorphic representation $\Psi[\phi]$, it is $L^2_{Hol}(\mathcal{N}'_\mathbb{C}, D\mu_c)$, for its respective complex Gaussian measure.\textit{Id est}, a pure state is to be thought as a function with domain in the DFN space $\mathcal{N}'$ or its complexification $\mathcal{N}'_\mathbb{C}$. For further details on the physical meaning of each representation we refer to the second part of this series  \cite{alonso2023a}.  

It is common knowledge that the Hilbert space of an scalar field theory is a bosonic Fock space. In this context, the notion of particle in QFT is defined as the eigenstates of the  particle number operator.   The precise relation between the Fock space and the $L^2$ Hilbert spaces is given by the Wiener-Ito decomposition in homogeneous chaos, which we discuss briefly below. In the same lines, we can also identify common tools of gaussian analysis as Malliavin derivatives and Skorokhod integrals with creation and annihilation operators defining in the process the notion of charge and particle number. 



\subsubsection{The Segal isomorphism: Wiener-Ito chaos decomposition and the Fock space}
The classical phase space $\mathcal{M}_C$ can be considered as the set of complex test functions $\mathcal{N}_\mathbb{C}$ once we introduce a suitable complex structure on the set $\mathcal{N}\times \mathcal{N}$ (containing the fields and their momenta). We will start dealing with the complex case and then particularize our result to the real case. If the Gel'fand triple discussed above is considered with respect to complex spaces, we can build the triple

\begin{equation}
\label{eq:complexgelfand}
\mathcal{M}_C\sim \mathcal{N}_\mathbb{C}\subset \mathcal{H}_\Delta\subset \mathcal{N}'_\mathbb{C},
\end{equation}
 where $\mathcal{H}_\Delta$ represents the, now complex, Cameron-Martin Hilbert space. By using this triple, we can consider $\mathcal{M}_C$ as a linear subspace of $\mathcal{N}'_\mathbb{C}$. This means that we will think at our classical fields and their momenta as distributions, dual to the nuclear space of complex smooth functions of compact support on $\Sigma$.

Now, we move to the level of second quantization. We therefore deal with the Hilbert space of quantum pure states which is in either case a subspace of $L^2(\mathcal{N}'_\mathbb{C}, D\mu_c)$. This  larger space is called the prequantum Hilbert space. For further insight in role in quantization see the second part of this series \cite{alonso2023a}.  In order to study its structure we will perform an orthogonal decomposition of the algebra of polynomials $\mathcal{P}(\mathcal{N}'_{\mathbb{C}})$ introduced above in  \eqref{eq:polynomialscompl}. We will see that this construction allows us to define an isomorphism with the (bosonic) Fock space constructed from the  (Cameron-Martin) Hilbert space $\mathcal{H}_\Delta$. Then we define  $\mathcal{W}_\mu^{:(n,m):}$ and denote it the space of homogeneous complex chaos of degrees $(n,m)$ as
\begin{equation}
\label{eq:wickspaces}
\mathcal{W}_\mu^{:(n,m):}=\overline{\mathcal{P}^{(n,m)}}\cap \overline{\mathcal{P}^{{(n-1,m)}}}^\perp\cap \overline{\mathcal{P}^{{(n,m-1)}}}^\perp,
\end{equation}
where $\overline{\mathcal{P}^{(n,m)}}$ represents the completion (in $L^2(\mathcal{N}'_\mathbb{C}, D\mu_c)$) of the set of polynomials of degrees $(n,m)$  and $\overline{\mathcal{P}^{(n-1,m)}}^\perp$ the orthogonal complement of the completion of the set of polynomials of degrees $(n-1,m)$. The sequence of spaces  $\mathcal{W}_\mu^{:(n,m):}$  contain a basis of orthogonal polynomials for the Hilbert space $L^2(\mathcal{N}'_\mathbb{C}, D\mu_c)$. In addition, these spaces are   generated by complex Wick monomials of the gaussian measure $\mu$ (or, equivalently, of the covariance $\Delta$)  generalizing the definition of complex Hermite polynomials to this more abstract setting \cite{Hida1980} as follows


\begin{definition}
    \label{def:ComplexWickMonomials}{\bf (Complex Wick monomials)} 
    Let  $\Delta^{\mathbf{xy}}$ be the covariance of the Gaussian measure $\mu$.  
    We will use $\mathcal{W}_\Delta(\phi^n\bar{\phi}^m)^{\vec{\mathbf{x}}\vec{\mathbf{y}}}$, with $\vec{x}=(x_i)_{i=1}^n, \vec{y}=(y_i)_{i=1}^m$, to refer to the distribution belonging to the symmetric tensor product of $n$ copies on $\mathcal{N}'_\mathbb{C}$ and the symmetric product of $m$ copies on ${\mathcal{N}'_\mathbb{C}}^*$ i.e. $\mathcal{W}_\Delta(\phi^n\bar{\phi}^m)^{\vec{\mathbf{x}}\vec{\mathbf{y}}}\in{\mathcal{N}'_\mathbb{C}}^{\hat{\otimes}n}\otimes {\mathcal{N}'_\mathbb{C}}^{*\hat{\otimes}m}$ such that
    \begin{equation}
    \label{eq:ComplexWienerIto}
    \mathcal{W}_\Delta(\phi^n\bar{\phi}^m)^{\vec{x}\vec{y}}:=\prod_{i=1}^n\prod_{j=1}^m\frac{\delta}{\delta \bar\rho_{x_i}}\frac{\delta}{\delta \rho_{y_j}}\exp\left[{\bar\rho_z\phi^z+\rho_z\bar\phi^z-\bar\rho_u\Delta^{uv}\rho_v}\right]\bigg\lvert_{\rho_{\mathbf{x}}=0}\;.
    \end{equation}
    We will refer to this distribution as \textit{complex Wick monomial} of degrees $(n,m)$ with respect to the covariance $\Delta^{\mathbf{xy}}$.
\end{definition}



 With this set of polynomials, we can write any state in the Hilbert space $L^2(\mathcal{N}_\mathbb{C}, D\mu_c) $ as a linear combination 
    \begin{equation}
    \label{eq:ComplexWienerItodec}
    \Psi(\phi^{\mathbf{x}},\bar\phi^\mathbf{y})=\sum_{n,m=0}^\infty \psi^{(n,\bar{m})}_{\vec{x}_n\vec{y}_m} \mathcal{W}_\Delta(\phi^n\bar{\phi}^m)^{\vec{x}_n\vec{y}_m}
    \end{equation}
    We call this decomposition Wiener-Ito chaos decomposition in homogeneous chaos.
    Notice that the set of polynomials can be decomposed in holomorphic and anti-holomorphic parts 
    \begin{equation}
    \label{eq:ComplexPolinomials}
    \mathcal{P}_{Hol}(\mathcal{N}'_{\mathbb{C}})\otimes \mathcal{P}_{\overline{Hol}}({\mathcal{N}_{\mathbb{C}}'}^*)\subset L^2\big({\mathcal{N}_{\mathbb{C}}'},D\mu_c\big)= L^2_{Hol}\big({\mathcal{N}_{\mathbb{C}}'}, D\mu_c\big)\bar{\otimes} L^2_{\overline{Hol}}\big({\mathcal{N}_{\mathbb{C}}'}, D\mu_c\big).
    \end{equation}
    This decomposition provides an exceptional tool for gaussian analysis, in terms of the coefficients on the Wiener-Ito decomposition we can write the scalar product as:
    \begin{align}
    \label{eq:complexchaosDecomposition}
    \int_{\mathcal{N}'_\mathbb{C}}D\mu_c(\phi^\mathbf{x})\overline{\Phi}(\phi^\mathbf{x})\Psi(\phi^\mathbf{x})
    &= \nonumber \\ \sum_{n,m=0}^\infty n!m!&\ \overline{\varphi^{(n,\bar{m})}}_{\vec{u}_n\vec{v}_m}\big[(\Delta^*)^n\big]^{\vec{u}_n\vec{x}_n} \big[\Delta^m\big]^{\vec{v}_m\vec{y}_m} \psi^{(n,\bar{m})}_{\vec{x}_n\vec{y}_m}
    =\nonumber \\ &
    \sum_{n,m=0}^\infty n!m! \big\langle\varphi^{(n,\bar{m})}_{\vec{\mathbf{x}}_n\vec{\mathbf{y}}_m},\psi^{(n,\bar{m})}_{\vec{\mathbf{x}}_n\vec{\mathbf{y}}_m}\big\rangle_{\mathcal{H}^{(n,\bar{m})}_\Delta}
    \end{align}

    Therefore, once weighted with $\sqrt{n!m!}$, the coefficients  of the Wiener-Ito decomposition belong to the Bosonic Fock space $\Gamma_{hol}\mathcal{H}_\Delta\otimes \Gamma_{\overline{hol}}\mathcal{H}_\Delta$.  To find the coefficients of the decomposition is useful to use the $S_{\mu}$ transform as
    \begin{equation}
    \label{eq:wienerItoSeagalwithTtransform}
    \psi^{(n,\bar{m})}_{\vec{x}_n\vec{y}_m}=\frac{1}{n!m!}\prod_{i=1}^n\prod_{j=1}^m K_{x_iu_i}K^*_{y_jv_j}\frac{\partial}{\partial \rho_{u_i}}\frac{\partial}{\partial \bar\rho_{v_j}}  {S_\mu[\Psi](\rho_\mathbf{x},\bar{\rho}_\mathbf{x})} \bigg\lvert_{\rho_\mathbf{x}=0}
    \end{equation}
    Let us recall that $K_{\mathbf{xy}}$ is the  (weak) inverse of the covariance $\Delta^{\mathbf{xy}}$.
     We see that the decomposition of \eqref{eq:complexchaosDecomposition} is such that
     \begin{equation}
        \label{eq:holantiholdec}
        L^2\big({\mathcal{N}_{\mathbb{C}}'}, D\mu_c\big)=L^2_{Hol}\big({\mathcal{N}_{\mathbb{C}}'}, D\mu_c\big)\bar{\otimes} L^2_{\overline{Hol}}\big({\mathcal{N}_{\mathbb{C}}'}, D\mu_c\big)=\sum_{n=0}^\infty\oplus \left( \sum_{m=0}^n\oplus \mathcal{H}_\Delta^{(n-m,\bar{m})} \right).
    \end{equation}

    This leads us to build an isomorphism between the Hilbert space and the Fock space which associates each state in $L^2({\mathcal{N}_{\mathbb{C}}'}, D\mu_c)$ and its coordinates in the linear decomposition:

\begin{definition}
    \label{def:SegalIso}{\bf (Segal isomorphism)} 
    We will call Segal isomorphism to the mapping


    \begin{equation}
\label{eq:Segal}
\begin{array}{cccc}
    \mathcal{I} :&  L^2\big({\mathcal{N}_{\mathbb{C}}'}, D\mu_c\big) & \longrightarrow &  \Gamma_{hol}\mathcal{H}_\Delta\otimes \Gamma_{\overline{hol}}\mathcal{H}_\Delta
 \\ 
 &  \Psi &\longmapsto &  \mathcal{I}(\Psi)=\left (\sqrt{n!m!}\psi^{(n,\bar{m})}_{\vec x_n \vec y_m} \right )_{n,m=0}^\infty 
\end{array}
\end{equation}
which defines a unitary isomorphism of Hilbert spaces considering the scalar product defined in \eqref{eq:complexchaosDecomposition}. In this decomposition the Fock space is further decomposed into the orthogonal direct sum $\Gamma_{hol}\mathcal{H}_\Delta\otimes \Gamma_{\overline{hol}}\mathcal{H}_\Delta
=\sum_{n=0}^\infty\oplus \left( \sum_{m=0}^n\oplus \mathcal{H}_\Delta^{(n-m,\bar{m})} \right)$.
\end{definition} 
\paragraph*{The holomorphic representation}
    The holomorphic representation is described by the holomorphic subspace $L^2_{Hol}\big({\mathcal{N}_{\mathbb{C}}'}, D\mu_c\big)$. Their Wick monomials are particularly simple $\mathcal{W}_\Delta(\phi^n\bar{\phi}^0)^{\vec{\mathbf{x}}}= (\phi^n)^{\vec{\mathbf{x}}}$. This fact simplifies remarkably computations in because holomorphic functions have a Wiener-Ito  decomposition given by 
    \begin{equation}
    \label{eq:holchaosdecomposition}
    \Psi_{Hol}(\phi^{\mathbf{x}})=\sum_{n=0}^\infty\psi^{n}_{\vec{x}_n}(\phi^n)^{\vec{x}_n}
    \end{equation}
     independently of the gaussian measure considered. Besides, the coefficients of each regular monomial belong to orthonormal subspaces. 

     \paragraph*{The Schrödinger representation}
    For the real case the subspace is $L^2\big({\mathcal{N}'}, D\mu\big)$. In this case we also have a Wiener-Ito decomposition.  We work out the decomposition from the real $S$ transform. Therefore Wick monomials are given by 


    \begin{equation}
    \label{eq:wickMonomialsGenerator}
    :\varphi^n:\lvert_\Delta^{\vec{\mathbf{x}}}=
    \prod_{i=1}^n\frac{\delta}{\delta \xi_{y_i}}
    \exp{\big(\xi_x\varphi^x-\frac12\xi_x\Delta^{xy}\xi_y\big)} \lvert_{\xi_{\mathbf{x}}=0}
    \end{equation}
    or recursively 
\begin{definition}
    \label{def:wickMonomials}{\bf (Wick monomials)} 
    We denote $:\varphi^n:\lvert_\Delta^{\vec{\mathbf{x}}}=\ :\varphi^n:\lvert_\Delta^{\mathbf{x_1,\cdots,x_n}}$, and call Wick monomial of degree $n$ with respect to the covariance $\Delta^{\mathbf{xy}}$, the distribution belonging to the symmetric tensor product of $n$ copies on $\mathcal{N}'$ i.e. $:\varphi^n:\lvert_\Delta^{\vec{\mathbf{x}}}\in(\mathcal{N}')^{n\hat{\otimes}}$ such that
    \begin{align}
        :\varphi^0:\lvert_\Delta&=1\nonumber\\
        :\varphi^1:\lvert_\Delta^{\mathbf{x}}&=\varphi^{\mathbf{x}}\nonumber\\
        :\varphi^n:\lvert_\Delta^{\vec{\mathbf{x}},{\mathbf{x}}_{n-1},{\mathbf{x}}_n}&=
        :\varphi^{n-1}:\lvert_\Delta^{\vec{{\mathbf{x}}},{\mathbf{x}}_{n-1}}\varphi^{{\mathbf{x}}_n}-(n-1)\Delta^{\mathbf{x}_n(\mathbf{x}_{n-1}}:\varphi^{n-2}:\lvert_\Delta^{\vec{\mathbf{x}})}
        \label{eqn:wickPowers}
    \end{align}
    where the parenthesized indices are symmetrized according to the convention  \eqref{eq:symmetrization}. 
\end{definition}
  Then the Wiener-Ito decomposition is given by
$$\Psi(\varphi^\mathbf{x})=\sum_{n=0}^\infty \psi^{(n)}_{\vec{x_n}}:\varphi^n:\lvert_\Delta^{\vec{{x}_n}}$$
 the coefficients belong to the weighted Bosonic Fock space obtained from $\Gamma\mathcal{H}_\Delta$ with weights $\sqrt{n!}$ and the scalar product is 
\begin{multline}
    \label{eq:complexchaosprod}
    \int_{\mathcal{N}}D\mu(\varphi^\mathbf{x})\overline{\Phi}(\varphi^\mathbf{x})\Psi(\varphi^\mathbf{x})
    = \\   \sum_{n=0}^\infty n!\ \overline{\phi^{(n)}}_{\vec{u}_n}\big[\Delta^n\big]^{\vec{u}_n\vec{x}_n}  \psi^{(n)}_{\vec{x}_n}
    = 
    \sum_{n,}^\infty n! \big\langle\phi^{(n)}_{\vec{\mathbf{x}}_n},\psi^{(n)}_{\vec{\mathbf{x}}_n}\big\rangle_{\mathcal{H}^{n}_\Delta}
    \end{multline}


    Notice that this recursive definition coincides with the usual prescription for normal ordering in quantum field theory where $\Delta^{\mathbf{xy}}$ is the propagator of the theory considered. This feature allows us to compute pointwise products of functions in terms of their chaos decompositions using the well known tool of Feynman diagrams \cite{Moshayedi2019}.


    \subsubsection{The Seagal-Bargmann tranform: Isomorphisms between the real and holomorphic cases.}
    In the case where the covariance for the complex case, $\Delta^{\mathbf{xy}}$ is real and symmetric we can stablish a unitary isomorphism between the subspace of holomorphic functions $L^2_{Hol}(\mathcal{N}'_{\mathbb{C}},D\mu_c)$ and  $L^2(\mathcal{N}',D\mu)$ just by performing an integration over the imaginary part of the coordinates:    


    \begin{definition}
        \label{def:Segaltranf}{\bf (Segal-Bargmann transform)} We denote by Segal-Bargmann transform the unitary isomorphism
        $$
        \begin{array}{cccc}
        \mathcal{B} :&  L^2(\mathcal{N}',D\mu)  & \longrightarrow &  L_{Hol}^2(\mathcal{N}_{\mathbb{C}}',D\mu_c) \\ 
         & \Psi_{Real} &\longmapsto & \Psi_{Hol}  
        \end{array}
        $$
        where both $\mu,\mu_c$ are gaussian measures with a real symmetric covariance $\Delta^{\mathbf{xy}}$, adapted in each case to its domain. To write the explicit formula for this isomorphism we follow a slightly different convention and write $\phi^\mathbf{x}=\varphi^\mathbf{x}+i\pi^\mathbf{x}$ to simplify computations. Then it follows   $\int_{\mathcal{N}'_\mathbb{C}} D\mu_c(\phi^{\mathbf{x}})=\int_{\mathcal{N}'} D\mu(\varphi^{\mathbf{x}})\int_{\mathcal{N}'} D\mu(\pi^{\mathbf{x}})$. Let $\Psi_{Hol}$ be a holomorphic function with chaos decomposition given by \eqref{eq:holchaosdecomposition}, integrating over the imaginary degrees of freedom 
    \begin{equation*}
        \label{eq:schHolrelation}
        \begin{split}
        \int D\mu(\pi^{\mathbf{x}})\Psi_{Hol}(\phi^\mathbf{x})=
        %\\ 
        \sum_{n=0}^\infty &\int D\mu(\pi^{\mathbf{x}}) \psi^{(n)}_{\vec{x}_n}{[(\varphi+%&
        i\pi)^n]^{\vec{x}_n}}=
        \\
        \sum_{n=0}^\infty  \psi^{(n)}_{\vec{x}_n}\frac{\partial}{\partial \xi^n_{\vec{x}_n}}\int D\mu(\pi^{\mathbf{x}})& e^{\xi_z(\varphi+i\pi)^z}\bigg\lvert_{\xi_{\mathbf{x}}=0}= 
        \\
        \sum_{n=0}^\infty  \psi^{(n)}_{\vec{x}_n}
        \frac{\partial}{\partial \xi^n_{\vec{x}_n}}\Bigg[&   e^{{\xi_z\varphi^z} } \int D\mu(\pi^{\mathbf{x}}) e^{i{\xi_z\pi^z}}\Bigg]_{\xi_{\mathbf{x}}=0}= 
        \\
        \hspace{ 2 em}\sum_{n=0}^\infty \psi^{(n)}_{\vec{x}_n}\frac{\partial^n}{\partial \xi^n_{\vec{x}_n}}e^{\xi_x\varphi^x-\frac{\xi_u\Delta^{uv}\xi_v}{2}}\Big\lvert_{\xi_\mathbf{x}=0}&=
        \sum_{n=0}^\infty \psi^{(n)}_{\vec{x}_n}:\varphi^n:\lvert_{{\Delta}}^{\vec{x}_n}=
         \Psi_{Real}(\varphi^{\mathbf{x}})
        \end{split}
    \end{equation*}
    where $\xi_\mathbf{x}\in \mathcal{N}$ is real.   With this consideration,  $\mathcal{B}(\Psi_{Real})$ is computed obtaining the chaos decomposition of  $\Psi_{Real}$ and replacing Wick monomials by regular holomorphic monomials. 
    
    A similar isomorphism is obtained using the antiholomorphic subspace in the analogous form
    
    $$
    \overline{\mathcal{B}} :L^2(\mathcal{N}',D\mu)   \longrightarrow  L_{\overline{Hol}}^2(\mathcal{N}_{\mathbb{C}}',D\mu_c).
    $$
\end{definition}


These transforms play a relevant role in the construction of physical quantum field theories, as they allow for the relation of the holomorphic, antiholomorphic and Schr\"odinger functional representations of quantum states with respect to their appropriate Gaussian measures.




\subsubsection{Malliavin calculus: creation and annihilation operators and charges}
\label{sssec:MaliiavinCalculus}
The space $L^2(\mathcal{N}',D\mu)$ is endowed with some useful properties that may be used to model further relevant tools in QFT such as creation and annihilation operators, or the number operator. Let us briefly introduce them in the simpler case of real fields and  generalize them later to the case of complex ones.

As in the complex case a polynomial of  a real variable $\varphi^{\mb{x}}\in \mathcal{N}'$,  with coefficients in $\mathcal{N}_{\mathbb{C}}$ and degree $n$  can be written as
\begin{equation}
\label{eq:realpolynomials} 
P_n(\varphi)= C^0+C^1_x\varphi^x+C^{2}_{x_1x_2}(\varphi^2)^{x_1x_2}+\cdots+C^n_{\vec{x}_n}(\varphi^n)^{\vec{x_n}}
\end{equation}    
 The algebra of polynomials of real variable is denoted by $\mathcal{P}(\mathcal{N}')$. This set is dense in $L^2(\mathcal{N}',D\mu)$ so we start introducing the operations of Malliavin calculus on $\mathcal{P}(\mathcal{N}')$ and then extend them to the whole space by linearity. On polynomials, the Fréchet derivative has a particularly useful expression in terms of the operator $\partial_{\varphi^\mathbf{y}}=\partial_{\mathbf{y}}$ that can be defined over monomials as 
\begin{equation}
\label{eq:Fréchetderivative}
\partial_{\varphi^y} C_{x_1,\cdots,x_n}(\varphi^n)^{x_1,\cdots,x_n}=nC_{y,x_1,\cdots,x_{n-1}}(\varphi^{n-1})^{x_1,\cdots,x_{n-1}} .
\end{equation}
By linearity, this operator is extended  to $\mathcal{P}(\mathcal{N}')$ without further complications.

The Malliavin derivative is the extension of the Fréchet derivative, defined in \eqref{eq:Fréchetderivative}, to an operator 
\begin{equation}
\label{eq:malliavinDerivative}
\partial_\mathbf{x}: \mathbb{D}_\mu^{2,1}\longrightarrow \mathcal{H}_\Delta\otimes L^2(\mathcal{N}',D\mu)
\end{equation}
where the domain of the operator is
\begin{equation}
    \label{eq:MalliavinSovoleb}
    \mathbb{D}_\mu^{2,1}=\left\{\sum_{n=0}^\infty\psi^{(n)}_{\vec{x}_n}:\varphi^n:\lvert_\Delta^{\vec{x}_n}\textrm{ s.t } \sum_{n=1}^\infty n\cdot n!\lVert\psi^{(n)}_{\vec{\mathbf{x}}_n}\rVert_{\Delta}^2< \infty \right\}
\end{equation}
and acts over $\Psi(\varphi^\mathbf{x})$ as

\begin{equation}
\label{eq:actionMalliavin}
\partial_x\Psi(\varphi^\mathbf{x})=\sum_{n=1}^\infty n\psi^{(n)}_{x,\vec{y}_{n-1}}:\varphi^{n-1}:\lvert_\Delta^{\vec{y}_{n-1}}
\end{equation}
The key feature of this derivative is that the resulting derivative makes sense as a $\mathcal{H}_\Delta$ function and does not have a pointwise meaning as  \eqref{eq:Fréchetderivative}. The Malliavin derivative can thus be considered as the generator of derivatives in directions of the Cameron-Martin Hilbert space. Furthermore, through the isomorphism $\mathcal{I}$, the contraction $\Delta^{xy}\partial_y$ is mapped to the annihilation operator on the Fock space $\Gamma(\mathcal{H}_\Delta)$.

The adjoint operator of the Malliavin derivative is the Skorokhod integral $\partial^{*\mathbf{x}}$. It is an operator defined over $\textrm{Dom}(\partial^{*\mathbf{x}})\subset\mathcal{H}_\Delta\otimes L^2(\mathcal{N}',D\mu)$ 

\begin{equation}
    \label{eq:SkorokhodIntegral}
    \begin{array}{cccc}
        \partial^{*\mathbf{x}}:& \textrm{Dom}(\partial^{*\mathbf{x}}) &\longrightarrow &L^2(\mathcal{N}',D\mu)\\
         & \Psi_{\mathbf{x}}(\varphi^{\mathbf{y}}) & \longmapsto & \partial^{*{x}} \Psi_{x}(\varphi^{\mathbf{y}})
    \end{array}
\end{equation}
It is important to remark that the Skorokhod integral is able to integrate more elements than those obtained from \eqref{eq:actionMalliavin}. Indeed, a general element $\Psi_{\mathbf{x}}\in\mathcal{H}_\Delta\otimes L^2(\mathcal{N}',D\mu)$ admits a Wiener-Ito decomposition in which 
\begin{equation}
\label{eq:parametricWienerIto}
\sum_{n=0}^\infty\psi^{(n)}_{x,\vec{y}_n}:\varphi^n:\lvert_\Delta^{\vec{y}_n}\hspace{0.25 cm} \textrm{ in general } \hspace{0.25 cm}  \psi^{(n)}_{(x,\vec{y}_n)}\neq \psi^{(n)}_{x,\vec{y}_n}
\end{equation}
In other words $(\psi^{(n)}_{x,\vec{y}_n})_{n=0}^\infty$ is not necessarily symmetrized and hence it may not be an  element of the Bosonic Fock space $\Gamma(\mathcal{H}_\Delta)$. Nonetheless if  
$\tilde{\psi}^{(n+1)}_{\vec{\mathbf{z}}_{n+1}}=\psi^{(n)}_{(\mathbf{x},\vec{\mathbf{y}}_n)}$ is indeed the decomposition of an element of $\Gamma(\mathcal{H}_\Delta)$ then $\Psi_{\mathbf{x}}\in  \textrm{Dom}(\partial^{*\mathbf{x}})$ and 

\begin{equation}
\label{eq:actionSkorokhod}
\partial^{*{x}}\Psi_{{x}}=\sum_{n=1}^{\infty} \tilde{\psi}^{(n)}_{\vec{z}_n}:\varphi^n:\lvert_\Delta^{\vec{z}_n}
\end{equation}
This, for a large class of functions in $\textrm{Dom}(\partial^{\star\mathbf{x}})$, can be written in terms of the Malliavin derivative as 
\begin{equation}
\label{eq:creationOperator}
\partial^{*{x}}=-\Delta^{xy}\partial_y+\varphi^x
\end{equation}

The Skorokhod integral is a generalization of the Ito integral, that coincides with it when the domain has dimension one that is identified with a time parameter, and we consider only $\Psi_x$ as a process adapted to the filtration of Brownian motion, see \cite{Nunno2009}.  This operator is mapped onto the creation operator on the Fock space picture by the isomorphism $\mathcal{I}$. 
It is even possible to extend these operations to the space of Hida distributions  $(\mathcal{N}_\mathbb{C})'$ introduced in the next section. For further details on Hida-Malliavin calculus,  we refer the interested reader to  \cite{Nunno2009}.

Another operator that can be easily written in terms of the latter two is the particle number 

\begin{equation}
\label{eq:particleNumber}
N_{\Delta}=\partial^{*x}\partial_x=-\Delta^{xy}\partial_x\partial_y+\varphi^x\partial_x
\end{equation}

These notions can be easily extended to the holomorphic and antiholomorphic cases, which will represent matter-antimatter  models. 
\begin{gather}
\label{eq:holomorphicMalliavin}
\partial_{\phi^x}\Psi(\phi^{\mathbf{x}},\bar\phi^\mathbf{y})=\sum_{n,m=0}^\infty n\mathcal{W}_\Delta(\phi^{n-1}\bar{\phi}^m)^{\vec{z}_{n-1}\vec{y}_m}\psi^{(n,\bar{m})}_{x,\vec{z}_{n-1}\vec{y}_m} \\
\partial_{\bar\phi^x}\Psi(\phi^{\mathbf{x}},\bar\phi^\mathbf{y})=\sum_{n,m=0}^\infty m\mathcal{W}_\Delta(\phi^{n}\bar{\phi}^{m-1})^{\vec{z}_{n}\vec{y}_{m-1}}\psi^{(n,\bar{m})}_{\vec{z}_{n}\vec{y}_{m-1},x} 
\end{gather}
Here each derivative must be understood in the Wirtinger sense 

\begin{equation}
\label{eq:wirtinger}
\partial_{\phi^x}=\frac{1}{\sqrt{2}}\big(\partial_{\varphi^x}-i \partial_{\pi^x} \big),\hspace{0.5 cm} \partial_{\bar\phi^x}=\frac{1}{\sqrt{2}}\big(\partial_{\varphi^x}+i \partial_{\pi^x} \big)
\end{equation}
On the other hand, the associated Skorokhod integrals act on (anti)holomorphic functionals as: 

\begin{align}
\label{eq:holomorphicSkorokhod}
\partial^{*\phi^x} =-\Delta^{xy}\partial_{\bar\phi^y}+\phi^x, & &\partial^{*\phi^x}\mathcal{W}_\Delta(\phi^{n}\bar{\phi}^m)^{\vec{z}_{n}\vec{y}_m}= \mathcal{W}_\Delta(\phi^{n+1}\bar{\phi}^m)^{x,\vec{z}_{n}\vec{y}_m}\nonumber \\
\partial^{*\bar\phi^x} =-(\Delta^*)^{xy}\partial_{\phi^y}+\bar\phi^x, & &\partial^{*\bar\phi^x}\mathcal{W}_\Delta(\phi^{n}\bar{\phi}^m)^{\vec{z}_{n}\vec{y}_m}= \mathcal{W}_\Delta(\phi^{n}\bar{\phi}^{m+1})^{\vec{z}_{n}\vec{y}_m,x}
\end{align}

In the same fashion as in the real case, we can write down the particle number operator $N$ and the $U(1)$ charge $Q$ encoding the matter/antimatter content of the states:
\begin{align}
\label{eq:ChargeAndNumber}
    N=\partial^{*\phi^x}\partial_{\phi^x}+ \partial^{*\bar\phi^x}\partial_{\bar\phi^x} & & Q=\partial^{*\phi^x}\partial_{\phi^x}- \partial^{*\bar\phi^x}\partial_{\bar\phi^x}
\end{align}

With those operators we see that the decomposition of \eqref{eq:complexchaosDecomposition} is such that

\begin{equation}
\label{eq:descompositionCharge}
L^2\big({\mathcal{N}_{\mathbb{C}}'},D\mu_c\big)=\sum_{n=0}^\infty\oplus \left( \sum_{m=0}^n\oplus \mathcal{H}_\Delta^{(n-m,\bar{m})} \right)
\end{equation}
where $\mathcal{H}_\Delta^{(n-m,\bar{m})}$ are simultaneous eigenspaces of $N$ and $Q$ with eigenvalues $n$ and $n-2m$ respectively. 


\subsection{Modelling quantum field pure states II: Hida test functions and Hida calculus}
\label{sec:Hida}
The set of Hida test functions  represents a {\it second quantized} notion of test functions that may be considered just an additional example of nuclear Fréchet space. Due to its relevance in the description of the space of quantum field pure states we have prepared a specific section for them and their tools. In order to do that, let us consider first the Hilbert space which we will use in the construction of a Gel'fand triple at this level. From its properties, we will be able to define the nuclear space of Hida test functions in a simple way. What we mean as second quantized construction is to consider that these Hida test functions are defined as  functions on a DFN space $\mathcal{N}'_{\mathbb{C}}$ with coefficients in its dual $\mathcal{N}_{\mathbb{C}}$. The topology of the latter is then provided by a sequence of Hilbert spaces $(\mathcal{H}_p, \langle \cdot, \cdot \rangle_p)$ and embeddings $\mathcal{H}_q\hookrightarrow \mathcal{H}_p$, which are Hilbert-Schmidt. We will use these ingredients to build another sequence of spaces and Hilbertian norms at the level of the functions on $\mathcal{N}'_{\mathbb{C}}$ and define, via a projective limit, the corresponding nuclear space structure for$(\mathcal{N}_{\mathbb{C}})$ see \eqref{eq:Hidatestfunction}.\footnote{Notation in this case, even though standard, could be misleading. Notice  that $(\mathcal{N}_{\mathbb{C}})$ represents Hida test functions with domain $\mathcal{N}'_{\mathbb{C}}$. } 




%\subsubsection{The nuclear space of Hida test functions and the second quantization Gel'fand triple}
 Let's consider the Gel'fand triple associated to the case when $\mathcal{N}_{\mathbb{C}}$ are complex test functions of compact support ({\it i.e.} the complex counterpart of equation \eqref{eq:riggedOnAManifold}) with respect to the Hilbert space $\mathcal{H}_{Vol}=L^2(\Sigma,\mathrm{dVol}_h)$. This Hilbert space is  obtained as the completion of $C^\infty_c(\Sigma)$ with respect to the scalar product induced by the measure $\mathrm{dVol}_h$ associated with a riemannian tensor $h$ on the hypersurface $\Sigma$. Thus we are considering the triple 
\begin{equation}
    \label{eq:riggedOnAManifoldcomp}
    \mathcal{N}_{\mathbb{C}}\subset L^2(\Sigma,\mathrm{dVol}_h) \subset \mathcal{N}_{\mathbb{C}}'
    \end{equation}

By using this rigged space, we can represent the space of test functions $C^\infty_c(\Sigma)_{\mathbb{C}}$ as a suitable linear subspace of the space of distributions $\mathcal{N}_{\mathbb{C}}\subset D'(\Sigma)_{\mathbb{C}}$. We shall use later this set as the set of classical fields.  Thus, our quantum pure states will be written then as test functions on this manifold $\mathcal{N}'_{\mathbb{C}}=D'(\Sigma)_{\mathbb{C}}$. Our notion of test function is provided by a nuclear topology and the injection of that topology into a Gel'fand triple respecting continuity.  This is the construction used in White Noise analysis \cite{Hida1993,Holden2010,Wester2003}. The topology of the space of test functions does not really depend on the particular Gel'fand triple that we use to realize the duality. For this reason is common to introduce them as a subset of the White Noise space of square integrable functions $L^2(\mathcal{N}'_{\mathbb{C}},D\beta)$ where $D\beta$ is the white noise whose characteristic functional is

\begin{equation}
\label{eq:whiteNoise}
\int_{\mathcal{N}'_{\mathbb{C}}}D\beta(\phi^\mathbf{x}) e^{i(\bar\rho_x\phi^x+\rho_x\bar\phi^x)}= e^{-\bar\rho_x\delta^{xy}\rho_y}=\exp\left(-\int_\Sigma \textrm{dVol}(x)\bar\rho(x)\rho(x)\right)
\end{equation}
Moreover, from this definition we have that $\mathcal{H}_{Vol}$ becomes the Cameron-Martin Hilbert space of $L^2(\mathcal{N}'_{\mathbb{C}},D\beta)$.


We will treat holomorphic and antiholomorphic functions separately. Thus we start by considering  holomorphic polynomials with coefficients in  $\mathcal{N}_{\mathbb{C}}$ as the minimal space of test functions in $L^2_{Hol}(\mathcal{N}'_{\mathbb{C}},D\beta)$. Recall that an holomorphic polynomial of degree $n$ with coefficients in $\mathcal{N}_{\mathbb{C}}$ can be written as
\begin{equation}
\label{eq:polynomials} 
P_n(\phi^{\mathbf{x}})= C^0+C^1_x\phi^x+C^{2}_{x_1x_2}(\phi^2)^{x_1x_2}+\cdots+C^n_{\vec{x}_n}(\phi^n)^{\vec{x_n}}
\end{equation}    
Consequently, we must consider  the algebra of holomorphic polynomials $\mathcal{P}_{Hol}(\mathcal{N}'_{\mathbb{C}})$, which we know that  form a dense subset of the Hilbert space $L^2_{Hol}(\mathcal{N}', D\beta)$. The next step is to enlarge the set $\mathcal{P}_{Hol}(\mathcal{N}'_{\mathbb{C}})$ to obtain a more operative notion of test function.   To that end we make use of the topology of $\mathcal{N}_{\mathbb{C}}$ to build a sequence of Hilbert spaces and injections, and define our final set as the projective limit of that sequence, as it was explained in \autoref{def:FN}.

 Remember that $\mathcal{N}_{\mathbb{C}}$ is a nuclear space that we see as a subset of $\mathcal{H}_{Vol}$, and therefore it admits the sequence of Hilbert spaces, scalar products and injections described in \autoref{def:FN}. If we consider the Fock space corresponding to the Hilbert space $\mathcal{H}_{Vol}$ from the Segal isomorphism, we can also consider a sequence of Fock spaces $\cdots \hookrightarrow \Gamma_{Hol} \mathcal{H}_p\hookrightarrow \cdots \hookrightarrow  \Gamma_{Hol} \mathcal{H}_{Vol}\cong L^2_{Hol}(\mathcal{N}'_{\mathbb{C}},D\beta)$, with Hilbertian norms $\| \cdot \|_p$ obtained from those of $\mathcal{H}_p$ (see \cite{Kondratiev1996}). The space of Hida test functions $(\mathcal{N}_{\mathbb{C}})$ is obtained as the projective limit of that sequence, and then we can write
\begin{equation}
\label{eq:Hidatestfunction}
(\mathcal{N}_{\mathbb{C}})=\bigcap_{p=0}^\infty \Gamma_{Hol} \mathcal{H}_p.
\end{equation}
It is important to remark that the topology of $(\mathcal{N}_{\mathbb{C}})$ depends only on the topology of the nuclear space $\mathcal{N}_{\mathbb{C}}$ (see \cite{Kondratiev1996}).

In an analogous way we can build the dual space $(\mathcal{N}_{\mathbb{C}})'$ as the injective limit
\begin{equation}
    \label{eq:Hidatestfunctionduak}
    (\mathcal{N}_{\mathbb{C}})'=\bigcup_{p=0}^\infty \Gamma_{Hol} \mathcal{H}_{-p}.
\end{equation}
With these spaces in hand, we can define the  White Noise  Gel'fand triple 
\begin{equation}
\label{eq:Triple_hida}
(\mathcal{N}_{\mathbb{C}})\subset L^2_{Hol}(\mathcal{N}'_{\mathbb{C}}, D\beta)\subset (\mathcal{N}_{\mathbb{C}})',
\end{equation}
  Notice that as $(\mathcal{N}_{\mathbb{C}})$ is a nuclear space, it is a convenient space to define a differential calculus on a manifold with it as a model and, with it, tensor fields and Hamiltonian structures. 

  As anticipated, the role of the white noise measure $D\beta$ is auxiliary in the construction. If we consider any other gaussian measure $D\mu_c$ with covariance $\Delta^{\mathbf{xy}}$ such that there exists a constant $C>0$ and a norm of the chain $\| \cdot \|_p$ such that $\lvert\bar\rho_x\Delta^{xy}\rho_y \rvert \leq C \| \rho_{\mathbf{x}} \|_p^2 $, then according to the characterization theorems \cite{Kondratiev1996} we also have 
  \begin{equation}
    \label{eq:Triple_hida_second}
    (\mathcal{N}_{\mathbb{C}})\subset L^2_{Hol}(\mathcal{N}'_{\mathbb{C}}, D\mu_c)\subset (\mathcal{N}_{\mathbb{C}})',
    \end{equation}
When the covariance of the measure encodes  some physical notion, as it may be the vacuum of the quantum field theory considered, this last triple would be the natural one to consider. We call this case \textbf{the second quantization Gel'fand triple}. We could as well present Hida test functions and distributions using this triple, but the fact that the topology of $(\mathcal{N}_{\mathbb{C}})$ only depends on that of $\mathcal{N}_{\mathbb{C}}$ turns  white noise into an auxiliary but simplifying way of presenting it.

Moreover, the vacuum of a theory is a particular feature of such theory and therefore if we want to consider every possible situation, it is better to model pure states as the common dense subspace of Hida test functions $(\mathcal{N}_{\mathbb{C}})$. In that case to recover the full space the only ingredient remaining is the gaussian measure that can be used to endow $(\mathcal{N}_{\mathbb{C}})$ with a pre-Hilbert topology and complete it to recover the whole space. This procedure is similar to the usual protocol for classical field theories which are modelled over spaces of test functions $\mathcal{N}$ and then completed with a Hilbertian norm to construct the quantum one particle space. We will use this fact in \cite{alonso2023a} to prescribe the dynamics of an arbitrary theory.  

 

We must also consider the decomposition for the antiholomorphic part. Thus, we consider the triple  
\begin{equation}
    \label{eq:gelfandtriples}
     (\mathcal{N}_\mathbb{C})\otimes (\mathcal{N}_\mathbb{C}^*) \subset L^2(\mathcal{N}'_{\mathbb{C}}, D\mu_c)\subset (\mathcal{N}_{\mathbb{C}})'\otimes (\mathcal{N}_\mathbb{C}^*)'.
    \end{equation}
     {Hence,} Hida test functions respect the decomposition in holomorphic and antiholomorphic parts of \eqref{eq:ComplexPolinomials}.
    


With this construction we have that $\mathcal{P}_{Hol}(\mathcal{N}'_{\mathbb{C}})\subsetneq (\mathcal{N}_{\mathbb{C}})$ is strict, as seen in the following example. 

\begin{example}{\bf (Hida-Malliavin calculus)}
    \label{ex:HidaMalliavin}
    The generator of the $S_\mu$ transform is a Hida test function
    $$\exp{\big({\rho}_x\bar\phi^x+{\bar\rho_x{\phi}^x}-\bar\rho_x\Delta^{xy}\rho_y\big)}\in (\mathcal{N}_\mathbb{C})\otimes (\mathcal{N}_\mathbb{C}^*)$$
    but it is not a polynomial. This is of key importance in order to extend $S_\mu[\Psi]$ to any Hida distribution $\Psi\in(\mathcal{N}_{\mathbb{C}})'\otimes (\mathcal{N}_\mathbb{C}^*)'$. In particular this means that we can find a chaos decomposition for any Hida distribution with a careful use of \eqref{eq:wienerItoSeagalwithTtransform}.  

    The generalization of the machinery of chaos decompositions to general distributions allows us to extend most of the tools of Malliavin calculus, described in \autoref{sssec:MaliiavinCalculus}, to distributions in $(\mathcal{N}_{\mathbb{C}})'\otimes (\mathcal{N}_\mathbb{C}^*)'$. This type of calculus is called Hida-Malliavin calculus and we refer elsewhere for a thorough construction \cite{Nunno2009,Holden2010}.
\end{example}



A common problem that is solved by Hida calculus is to make sense out of Radon-Nikodym derivatives. This will be particularly relevant when QFT theories with respect to evolving backgrounds are considered since these properties of Hida calculus will provide us with tools to compare the states of the quantum fields with respect to them. We can define 

\begin{example}{\bf (Translation of a gaussian measure)}
    \label{ex:translations} A general result of gaussian analysis, known as Cameron-Martin theorem, states that,  a translation in the domain by a constant element $\chi^\mathbf{x}\in \mathcal{N}'_\mathbb{C}$ provides a gaussian measure mutually singular with the original  one, unless $\chi^\mathbf{x}\in \mathcal{H}_K$, the Cameron-Martin Hilbert space. This implies that the \textit{would be Radon-Nikodym derivative}
    $$\frac{D\mu_c(\phi^\mathbf{x}-\chi^\mathbf{x})}{D\mu_c(\phi^\mathbf{x})}=\exp{\big(\bar\phi^xK_{xy}\chi^y+\bar\chi^xK_{xy}\phi^y-\bar\chi^xK_{xy}\chi^y\big)}$$
does not exist as an $L^1(\mathcal{N}'_\mathbb{C},D\mu_c)$ function. The power of Hida calculus is to extend its meaning to general $\chi^\mathbf{x}\in \mathcal{N}'_\mathbb{C}$ as an element of $(\mathcal{N}_{\mathbb{C}})'\otimes (\mathcal{N}_\mathbb{C}^*)'$ and provide a tool to perform  calculus with the generalized Wiener-Ito decomposition. 
\end{example}
 

\begin{example}{\bf (Mean zero gaussian measure)}
\label{ex:meanZeroGaussian} In general, gaussian measures with different covariances, even if they both fulfill an estimate of the type $\lvert\rho_x\Delta^{xy}\rho_y \rvert \leq C \| \bar\rho_{\mathbf{x}} \|_p^2 $, are mutually singular.  

 When working with the space of Hida test functions, if $\nu_c$ is
 another mean zero gaussian measure with covariance $\Omega^{\mathbf{xy}}$ then 
$$S_{\mu_c}\Big[\frac{D\nu_c}{D\mu_c}\Big](\rho_\mathbf{x})=e^{\bar\rho_x(\Omega^{xy},-\Delta^{xy})\rho_y},$$
where $\frac{D\nu_c}{D\mu_c}\in (\mathcal{N}_{\mathbb{C}})'\otimes (\mathcal{N}_\mathbb{C}^*)'$ makes sense as a Hida distribution even if it does not  as a function.  
\end{example}


\begin{example}{\bf (Reproducing Kernel)}
    \label{ex:repKernel} Coherent states of the holomorphic subspace $\mathcal{K}_{\chi}(\phi^{\mathbf{y}})\in L^2_{Hol}\big({\mathcal{N}_{\mathbb{C}}'},D\mu_c\big)$ can be extended to elements $\rho^{\mathbf{x}}\in \mathcal{H}_\Delta$ as $ \mathcal{K}_\rho(\phi^{\mathbf{y}})=e^{\bar\rho^xK_{xy}\phi^x}$ that makes sense as an element of  $L^2_{Hol}\big({\mathcal{N}_{\mathbb{C}}'},D\mu_c\big)$ but, moreover 
    $$S_{\mu_c}\Big[ e^{\bar\varphi^xK_{xy}\phi^x}\Big](\rho_\mathbf{x})=e^{\bar\varphi^x\rho_x}$$
    Then  $e^{\bar\varphi^xK_{xy}\phi^x}\in (\mathcal{N}_{\mathbb{C}})'$ 
    is a Hida distribution and works as a reproducing kernel when it is restricted to $(\mathcal{N}_{\mathbb{C}})\subset L^2_{Hol}\big({\mathcal{N}_{\mathbb{C}}'},D\mu_c\big)$. This is for $\Psi\in  (\mathcal{N}_{\mathbb{C}})$
    \begin{equation}
    \label{eq:rpekernelHida}
    \int_{\mathcal{N}_{\mathbb{C}}'}D\mu_c(\phi^{\mathbf{x}})e^{\bar\phi^xK_{xy}\varphi^x}\Psi(\phi^{\mathbf{x}})= \Psi(\varphi^x)
    \end{equation}
    \end{example}
\noindent See \cite{Hida1993} for further details and examples. 

\section{Adapting Classical Field Theory to quantization}
\label{sec:CFTQuantizar}
Before considering the set of quantum fields in the second part \cite{alonso2023a}, our first task is to adapt to quantization the geometrical description of the classical field theory that we are going to quantize. Alternative geometrical descriptions of classical field theories in terms of jet bundles also exist and admit  well defined Hamiltonian formalisms. Nonetheless, our goal is to build, from this classical manifold, a geometrical description of the space of quantum states on the space of fields, and our choice provides us with a simpler analog of the finite dimensional construction for quantum mechanical systems (see \cite{Kibble1979}).

We will focus on the real scalar field, i.e. the phase space of the theory will be given by the cotangent bundle to a real line bundle $\pi_{\Sigma L}:L\to \Sigma$. For simplicity we will consider the trivial case in which $L=\mathcal{N}$ is a space of smooth functions over $\Sigma$.

Furthermore, the set of fields modelled by any of these function spaces should be solutions of a certain dynamical field equation representing the type of field. The simplest example will be the case of a Klein-Gordon field theory. In this work we will consider only linear equations and therefore the corresponding manifold of fields will admit a linear structure which will significantly simplify the geometrical setting. 

\subsection{Choosing the space of classical fields and their momenta}

If $\mathcal{N}$ represents the set of fields, the corresponding phase space should correspond to the cotangent bundle $T^*\mathcal{N}\simeq \mathcal{N}\times \mathcal{N}'$, where we used the linearity of $\mathcal{N}$ and the corresponding triviality of the cotangent bundle. The coordinate representation of an element of this space is given by $(\varphi_\mathbf{x},\pi^\mathbf{y})$, with the coordinate systems described above. This space is endowed with a weakly symplectic form  which takes the local expression:

\begin{equation}
\label{eq:symplecticForm}
\omega_{T^*\mathcal{N}}=  d\pi_x \wedge d\varphi^x.
\end{equation}

This situation is problematic in the sense that the model space for the fiber and the base of the cotangent bundle $T^*\Gamma(\Sigma,L)$ are different. To quantize the system is convenient to express both coordinates over the same model space. To do so we endow $\Sigma$ with a Riemannian structure $h$, use it to define a measure $d\mathrm{Vol}$ and use the natural structure of rigged Hilbert space given by 
\begin{equation}
\label{eq:rhsOnSigma}
\mathcal{N}\subset L^2(\Sigma,d\textrm{Vol})\subset \mathcal{N}'
\end{equation} 

As a result, the elements of $\mathcal{N}$ corresponds to a linear subspace of its dual and we can write the coordinate representation of the image of an element of $\mathcal{N}$ as $\delta^{xy}\varphi_x$. Thus, in this setting, the manifold of fields must  be $\mathcal{M}_F=\mathcal{N}'\times \mathcal{N}'$ and coordinates $(\varphi^\mathbf{x},\pi^\mathbf{y})$ are to be understood as Darboux coordinates for a densely defined weakly symplectic structure
\begin{equation}
\label{eq:symplecticFormlifted}
\omega_{\mathcal{M}_F}=\delta_{xy} d\pi^x \wedge d\varphi^y =\int_\Sigma \frac{d^dx}{\sqrt{\lvert h\rvert}}\ [d\pi(x)\otimes d\varphi(x)-d\varphi(x)\otimes d\pi(x)],
\end{equation}
where now $\pi(x),\varphi(x)$ are densities of weight 1. Notice that this symplectic structure is defined over the dense subspace $L^2(\Sigma,d\textrm{Vol})\subset \mathcal{N}'$. If we consider the space of polynomial functions  $\mathcal{P}(\varphi^\mathbf{x},\pi^\mathbf{y})$ with coefficients in $\mathcal{N}$ ( that, eventually, can be  extended to more general completions) we can also consider a Poisson bivector of the form:
\begin{multline}
\label{eq:poissonbracket}
\left\{P,Q\right\}_{\mathcal{M}_F}=
\delta^{xy}\left( \partial_{\varphi^x}P\partial_{\pi^y}Q-\partial_{\pi^x}P\partial_{\varphi^y}Q \right):=\\
\int_\Sigma d\textrm{Vol}(x)\left( \frac{\partial P}{\partial\varphi(x)}\frac{\partial Q}{\partial\pi(x)}-\frac{\partial P}{\partial\pi(x)}\frac{\partial Q}{\partial\varphi(x)} \right),
\end{multline}
where partial derivatives are understood as Fréchet derivatives, which are well defined for $\mathcal{P}(\varphi^\mathbf{x},\pi^\mathbf{y})$. This Poisson structure is often mistaken with a symplectic structure in QFT quantization schemes since, for linear spaces, the difference is subtle see \cite{Brunetti2015} chapter 5 for further discussion on that matter. With this Poisson bracket at hand we can proceed as usual and write down Hamiltonian dynamics of classical observables.  Classical observables are elements of a suitable completion of the set of polynomials in the field and momenta variables $\mathcal{P}(\varphi^\mathbf{x},\pi^\mathbf{y})$ and existence and uniqueness of such dynamics is, in general, strongly dependent on this completion.   Such problems, though, will not be covered in this work.

Summarizing: our classical phase space shall be considered to be the linear subspace of $\mathcal{M}_F=\mathcal{N}'\times \mathcal{N}'$ corresponding to the injection of $\mathcal{N}$ by the rigged Hilbert space structure introduced on $L^2(\Sigma, d\textrm{Vol})$ by the riemannian structure of the Cauchy hypersurface $\Sigma$. On this space we can consider two (almost) equivalent structures: a (weakly) symplectic structure $\omega_{\mathcal{M}_F}$ and a Poisson tensor $\{ \cdot, \cdot \}_{\mathcal{M}_F}$.  

It is important to remark a subtle aspect of the construction: while the symplectic form on $T^*\mathcal{N}$ (Equation \eqref{eq:symplecticForm}) is canonical, the one on $\mathcal{M}_C$ (Equation \eqref{eq:symplecticFormlifted}   is not, since it depends on the choice of measure $d\mathrm{Vol}$ which we used to define the Gel'fand triple. The same is true for the Poisson bracket \eqref{eq:poissonbracket}. This aspect will be important when considering quantum classical hybrid models.


Now, we will learn to introduce another structure on it to define more conveniently the quantization of the classical model. In order to do that, we need to introduce a complex structure on $\mathcal{M}_F$.

\subsection{Introducing a complex structure}
In order to obtain a system better  suited for geometric quantization we need an additional (almost) complex structure $J_{\mathcal{M}_F}$ satisfying:
\begin{itemize}
	\item it is densely defined,
	\item  it leaves $D(J_{\mathcal{M}_F})\subseteq L^2(\Sigma,d\textrm{Vol})\subset \mathcal{N}'$ invariant
	\item and such that $J_{\mathcal{M}_F}^2=-\mathds{1}$ in its domain $D(J_{\mathcal{M}_F})$.
\end{itemize}

We shall see  in \cite{alonso2023a} that the complex structure required for quantization will be determined by the classical Hamiltonian on the space of fields. Hence, in order to consider different possible classical dynamics, we must explore the space of possible complex structures on $\mathcal{M}_F$.  It is simple to verify that there are infinitely many different structures satisfying those conditions. To express the most general complex structure let $(\varphi^{\mathbf{x}},\pi^{\mathbf{y}})$ be Darboux coordinates for $\mathcal{M}_F$ and let us consider the matrix elements $(A^t)^{\mathbf{x}}_\mathbf{y}:= \delta_{u \mathbf{y}} A^u_v\delta^{v\mathbf{x}}$. Notice that the definition of these variables already depends on the choice of the Hilbert space structure for the Gel'fand triple.

An arbitrary (almost) complex structure on $\mathcal{M}_F$ is locally expressed as 
\begin{multline}
\label{eq:complexStructure}
-J_{\mathcal{M}_F} =
(\partial_{\varphi^y},\partial_{\pi^y})\left( \begin{array}{cc}
A^y_x & \Delta^y_x \\
D^y_x & -(A^t)^y_x
\end{array} \right)  \left( \begin{array}{c}
d\varphi^{x} \\
d\pi^{x}
\end{array} \right)=\\
=  \partial_{\varphi^y} \otimes[A^y_xd\varphi^x+\Delta^y_xd\pi^x]+ \partial_{\pi^x}\otimes [D^y_xd\varphi^x-(A^t)^y_xd\pi^y]
\end{multline}
With $(A\Delta)^x_y=A^x_z\Delta^z_y$ we have $A^2+\Delta D=-\mathds{1},\ \Delta^t=\Delta, D^t=D, A\Delta=\Delta A^t$ and $A^tD=DA$. Notice that $D$ is fixed once $\Delta$  and $A$ are known. Let $K^{\mathbf{x}}_{\mathbf{y}}$ be the inverse of $\Delta_\mathbf{y}^{\mathbf{x}}$, i.e. $\Delta^{x}_zK^z_{y}=\delta^x_y$ then $D=(iA^t+\mathds{1})K(iA-\mathds{1})$.

 

This complex structure, together with the symplectic form \eqref{eq:symplecticFormlifted} induces a (pseudo)-Riemannian structure $\mu_{J,\mathcal{M}_F}(\cdot,\cdot)=\omega_{\mathcal{M}_F}(\cdot,-J\cdot)$, which, in the coordinates above, reads: 
\begin{multline}
\label{eq:riemannianInduced}
\mu_{\mathcal{M}_F}=(d\varphi^{x},d\pi^{x}) \left( \begin{array}{cc}
    -D_{xy} & A^t_{xy} \\
    A_{xy} & \Delta_{xy}
    \end{array} \right) \left( \begin{array}{c}
        d\varphi^{y} \\
        d\pi^{y}
    \end{array} \right)=\\
    -D_{xy}d\varphi^x\otimes d\varphi^y+\Delta_{xy}d\pi^x\otimes d\pi^y+ A_{xy}(d\varphi^y\otimes d\pi^x+d\pi^x\otimes d\varphi^y)
\end{multline}
where $A_{xy}=\delta_{xz}A^z_y$. To obtain a Riemannian structure form here we must impose also $\Delta_{\mathbf{xy}}>0>D_{\mathbf{xy}}$. 

With a change of coordinates, locally given by
\begin{align}
\label{eq:changeOfCoordinates}
d{\tilde{\varphi}^x}&=d{\varphi^x} &
d{\tilde{\pi}^x}&=A^x_yd{\varphi^y}+\Delta^x_yd{\pi^y} 
\\
\partial_{\tilde{\varphi}^x}&=\partial_{\varphi^x}- (K A)^y_x\partial_{\pi^y} &
\partial_{\tilde{\pi}^x}&=K^y_x\partial_{{\pi}^y}
\end{align}
 we can put in a canonical way $J_{\mathcal{M}_F}=\partial_{\tilde{\pi}^x}\otimes d\tilde{\varphi}^x-\partial_{\tilde{\varphi}^x}\otimes d\tilde{\pi}^x$, $\mu_{\mathcal{M}_F}=K_{xy}(d\tilde{\varphi}^x\otimes d\tilde{\varphi}^y+d\tilde{\pi}^x\otimes d\tilde{\pi}^y)$ and $\omega_{\mathcal{M}_F}=K_{xy} d\tilde{\pi}^x\wedge d\tilde{\varphi}^y$. Notice that we can cast the structures into this canonical form with other changes of variables, but we choose this one because it leaves $\varphi^\mathbf{x}$ (the manifold of fields) invariant. Notice that we are transforming the original cotangent bundle structure we began with as the new coordinates combine base and fiber coordinates of the original bundle structure.

Having chosen a complex structure, we can define an isomorphism for $\mathcal{M}_F$ to become a complex vector space $\mathcal{M}_C=\mathcal{M}_F^{\mathbb{C}}$. Notice that while $\mathcal{M}_F$ depends on the choice of the Hilbert space structure for the Gel'fand triple (and the original $T^*\mathcal{N}$), the definition of $\mathcal{M}_C$ depends also on the choice of the complex structure.  This fact will be important later. Notice, anyway, that the symplectic structures on $T^*\mathcal{N}'$, $\mathcal{M}_F$ and $\mathcal{M}_C$ are all diffeomorphic, but the diffeomorphisms depend on the metric structure $h$ on $\Sigma$ (for $\mathcal{M}_F$)  and the complex structure $J_{\mathcal{M}_C}$ (for $\mathcal{M}_C)$. 

 Having chosen the complex manifold $\mathcal{M}_C$, we consider holomorphic and antiholomorphic  coordinates which are locally written as $(\phi^{\mathbf{x}},\bar{\phi}^{\mathbf{y}})=\frac{1}{\sqrt{2}}(\tilde{\varphi}^{\mathbf{x}}+i\tilde{\pi}^{\mathbf{x}},\tilde{\varphi}^{\mathbf{y}}-i\tilde{\pi}^{\mathbf{y}})$. In this new manifold the complex structure becomes  $J_{\mathcal{M}_C}$, which  is now written as: 

\begin{equation}
\label{eq:complexholom}
J_{\mathcal{M}_C}=i \partial_{\phi^x}\otimes d\phi^x-i \partial_{\bar{\phi}^x}\otimes d\bar{\phi}^x.
\end{equation}
At the tangent space level we can always write the projection operators $P_{\pm}=\frac12(\mathds{1}\mp iJ)$ such that $P_{+}\partial_{\phi^{\mathbf{x}}}=\partial_{\phi^{\mathbf{x}}}$ and $P_{-}\partial_{\bar\phi^{\mathbf{x}}}=\partial_{\bar\phi^{\mathbf{x}}}$. We shall use these projections in the following sections.

As the triple $(\mu_C, \omega_C, J_C)$ defines a Kähler structure on $\mathcal{M}_C$, there exists
a (non unique) globally defined on a dense subspace (because the manifold is linear) Kähler potential  $\mathcal{K}$

\begin{equation}
\label{eq:kahlerpotential}
\mathcal{K}(\bar \phi^x, \phi^x)=\bar \phi^x  K_{xy} \phi^y .
\end{equation}



From the potential we can define the corresponding  Kähler form:

\begin{equation}
\label{eq:kahlerform}
\omega_{\mathcal{M}_C}= i\partial \bar \partial \mathcal{K},
\end{equation}
where $\partial$ and $\bar \partial$ represent the Dolbeault operators. It is important to notice that, as usually happens in the literature\footnote{See for example \cite{Woodhouse1997}}, the convention on the sign change from real to holomorphic coordinates, this is $\omega_{\mathcal{M}_C}=-\omega_{\mathcal{M}_F}$. With the Kähler form and the expression of $J_{\mathcal{M}_C}$ of \eqref{eq:complexholom} we obtain the expression of the riemannian form $\mu_{\mathcal{M}_C}$ and the hermitian form: $$h_{\mathcal{M}_C}= \frac{\mu_{\mathcal{M}_C}+i\omega_{\mathcal{M}_C}}{2}= K_{xy}d\phi^x\otimes d\bar{\phi}^y,$$ 
with inverse $h^{-1}_{\mathcal{M}_C}=\Delta^{xy}\partial_{\phi^x}\otimes\partial_{\bar\phi^y}$.

Finally, the expression of the Poisson bracket \eqref{eq:poissonbracket} in holomorphic coordinates becomes:
\begin{equation}
\label{eq:poissonHol}
\left\{P,Q\right\}_{\mathcal{M}_C}=
i\Delta^{xy}\left( \partial_{\phi^x}P\partial_{\bar\phi^y}Q-\partial_{\bar\phi^x}P\partial_{\phi^y}Q \right).
\end{equation}


It is important to notice that the election of the rigged Hilbert space made in \eqref{eq:rhsOnSigma}, to express the system in Darboux coordinates, is different for holomorphic coordinates and the sign of \eqref{eq:poissonbracket} is reversed due to our conventions. In the first case, we built the triple from the real classical fields onto the real momentum fields. In holomorphic coordinates the natural triple is defined from the complexified classical fields 
$\mathcal{N}_{\mathbb{C}}$ onto the complexified momentum fields $\mathcal{N}_{\mathbb{C}}'$:

\begin{equation}
\label{eq:HolomorphicTriple}
\mathcal{N}_{\mathbb{C}}\subset \mathcal{H}_\Delta\cong \mathcal{H}_K^*\subset {\mathcal{N}_{\mathbb{C}}'}
\end{equation}
 where $\mathcal{H}_\Delta=\overline{(\mathcal{N}_{\mathbb{C}},\langle \cdot,\cdot\rangle_\Delta)}$ is the  Hilbert space obtained by completion of $\mathcal{N}_{\mathbb{C}}$ with the scalar product induced by  the tensor $\Delta$ as $\langle \chi_{\mathbf{x}},\xi_{\mathbf{y}}\rangle_\Delta=\bar{\chi}_x\Delta^{xy}\xi_y$.  This space is identified with the dense subset of $\mathcal{N}_{\mathbb{C}}'$ on which the hermitian form $h_{\mathcal{M}_C}$ is finite. This Hilbert space  has a physical interpretation for $\mathcal{H}_\Delta$ as the space corresponding to the \textit{one-particle-state} structure. The space comprises the initial conditions for all the solutions of 'positive frequency' that are to be interpreted as particles in contrast with the elements of the conjugate $\mathcal{H}_{\Delta}^*$ which represent the corresponding anti-particles.    

On the other hand, we can consider a similar construction from the dual perspective if we consider the completion of  the dense subset of $\mathcal{N}_{\mathbb{C}}'$ on which the scalar product $\langle \cdot,\cdot\rangle_K$ induced by the hermitian tensor $K$ is defined.  Both spaces $\mathcal{H}_\Delta$ and $\mathcal{H}_K$ are identified by means of the Riesz representation theorem,  realizing the duality  $(\mathcal{N}_{\mathbb{C}},\mathcal{N}_{\mathbb{C}}')$ with the scalar product  $\langle \cdot,\cdot\rangle_\Delta$.  In the following, we will consider them as the same space.

Furthermore, the elements of $\mathcal{H}_K$ will be important later when we build the second quantization of our field theory. Indeed, when considering "wave-functions" on our manifold of classical fields, $\mathcal{H}_K$ is identified with the Cameron-Martin space represent the allowed translations of the unique Gaussian measure which has covariance $\Delta$. From that point of view, we will consider them to be the natural representation of the tangent space to our space of fields.

Summarizing: the classical phase space of fields is described as a Kähler complex manifold $\mathcal{M}_C$, which has associated a rigged Hilbert space structure encoding the physical \textit{one-particle-space} structure. In the second part of this series \cite{alonso2023a}, we will use these structures to define, by geometric quantization, a geometrical description of the quantum field theory corresponding to this  classical phase space of fields. We will also investigate, using the tools presented so far, different methods of quantization and the relations between them.


\section{Conclusions}
\label{sec:conclusions}
 

In this paper we have summarized and collected some important results of gaussian analysis with a vision in quantum field theory. First, we have argued that the natural domain for integration are strong duals of Nuclear Frechèt spaces (DFN) denoted by $\mathcal{N}'$ and that the concept of rigged Hilbert space emerged from the necessity of identify the Frechèt Nuclear (FN) spaces with subsets of their duals. 

 Then we defined a gaussian measure $D\mu_c$ and we have studied gaussian integration theory of functions with domain in the model DFN space. The space of functions analyzed is $L^2\big({\mathcal{N}_{\mathbb{C}}'}, D\mu_c\big)$ and we distinguished three classes of dense subsets of functions  that unveil different properties of the space that will be used in the second part of this series \cite{alonso2023a} to study different representations of a QFT. 
 
 The first example of dense subset are trigonometric exponentials $\mathcal{T}(\mathcal{N}'_\mathbb{C})$ that serve as a basis for an isomorphism with a reproducing kernel Hilbert space. This tool will be of huge importance in the second part of this series to stablish the properties of the algebra of observables. Closely related is the algebra of coherent states $\mathcal{C}_{Hol}(\mathcal{N}'_\mathbb{C})$ that is dense in the subspace of square integrable holomorphic functions $L^2_{Hol}\big({\mathcal{N}_{\mathbb{C}}'}, D\mu_c\big)$ and helps to stablish the analytical properties of that subspace as well as the decomposition  in holomorphic and antiholomorphic parts of the whole space.

 The last space of dense subset of $L^2\big({\mathcal{N}_{\mathbb{C}}'}, D\mu_c\big)$ are polynomials  $\mathcal{P}(\mathcal{N}'_\mathbb{C})$. This set helps to stablish tools of huge importance for gaussian analysis an its relation with quantum field theory. First, we established the Wiener-Ito decomposition theorem that provides an isomorphism with a bosonic Fock space completely determined by the properties of the gaussian measure. This will yield the particle interpretation of the QFT in \cite{alonso2023a}. Secondly, we developed the Segal-Bargmann transform that relates $L^2_{Hol}\big({\mathcal{N}_{\mathbb{C}}'}, D\mu_c\big)$ with its counterpart with real domain $L^2\big({\mathcal{N}'}, D\mu\big)$ and will serve the basis to relate the Holomorphic and Schrödinger representations in the second part of this series. Lastly we introduced the tools of integro-differential calculus, namely the Malliavin derivative and Skorokhod Integral that will represent creation and annihilation operators in the QFT.  



 Finally have described how the choice of Hida test functions to model our quantum states provides us with several tools which will be proven useful in the second part of this paper. Nonetheless, in order to be able to use them in our quantum field model, it is necessary to exhibit in some detail the origin of the geometric elements which is required to incorporate to the corresponding classical model. This is what we introduced in the last section of the present chapter, with particular emphasis in the Kähler structure which will determine the quantization of the theory. We will consider the implications of these objects in the second part of the paper.

 

\textbf{Acknowledgments}
The authors would like to thank  Carlos Escudero Liébana for pointing out to very useful bibliography on the connection of Malliavin calculus with QFT.
The authors acknowledge partial financial support of Grant PID2021-
123251NB-I00 funded by MCIN/AEI/10.13039/501100011033 and by the
European Union, and of Grant E48-23R funded by Government of Aragón.
C.B-M and D.M-C acknowledge financial support by Gobierno de Aragón
through the grants defined in ORDEN IIU/1408/2018 and ORDEN
CUS/581/2020 respectively.




%\addcontentsline{toc}{chapter}{Bibliography}
 

%\bibliography{Geom}
%\end{document}

\begin{thebibliography}{26}
    % BibTex style file: bmc-mathphys.bst (version 2.1), 2014-07-24
    \ifx \bisbn   \undefined \def \bisbn  #1{ISBN #1}\fi
    \ifx \binits  \undefined \def \binits#1{#1}\fi
    \ifx \bauthor  \undefined \def \bauthor#1{#1}\fi
    \ifx \batitle  \undefined \def \batitle#1{#1}\fi
    \ifx \bjtitle  \undefined \def \bjtitle#1{#1}\fi
    \ifx \bvolume  \undefined \def \bvolume#1{\textbf{#1}}\fi
    \ifx \byear  \undefined \def \byear#1{#1}\fi
    \ifx \bissue  \undefined \def \bissue#1{#1}\fi
    \ifx \bfpage  \undefined \def \bfpage#1{#1}\fi
    \ifx \blpage  \undefined \def \blpage #1{#1}\fi
    \ifx \burl  \undefined \def \burl#1{\textsf{#1}}\fi
    \ifx \doiurl  \undefined \def \doiurl#1{\url{https://doi.org/#1}}\fi
    \ifx \betal  \undefined \def \betal{\textit{et al.}}\fi
    \ifx \binstitute  \undefined \def \binstitute#1{#1}\fi
    \ifx \binstitutionaled  \undefined \def \binstitutionaled#1{#1}\fi
    \ifx \bctitle  \undefined \def \bctitle#1{#1}\fi
    \ifx \beditor  \undefined \def \beditor#1{#1}\fi
    \ifx \bpublisher  \undefined \def \bpublisher#1{#1}\fi
    \ifx \bbtitle  \undefined \def \bbtitle#1{#1}\fi
    \ifx \bedition  \undefined \def \bedition#1{#1}\fi
    \ifx \bseriesno  \undefined \def \bseriesno#1{#1}\fi
    \ifx \blocation  \undefined \def \blocation#1{#1}\fi
    \ifx \bsertitle  \undefined \def \bsertitle#1{#1}\fi
    \ifx \bsnm \undefined \def \bsnm#1{#1}\fi
    \ifx \bsuffix \undefined \def \bsuffix#1{#1}\fi
    \ifx \bparticle \undefined \def \bparticle#1{#1}\fi
    \ifx \barticle \undefined \def \barticle#1{#1}\fi
    \ifx \bconfdate \undefined \def \bconfdate #1{#1}\fi
    \ifx \botherref \undefined \def \botherref #1{#1}\fi
    \ifx \url \undefined \def \url#1{\textsf{#1}}\fi
    \ifx \bchapter \undefined \def \bchapter#1{#1}\fi
    \ifx \bbook \undefined \def \bbook#1{#1}\fi
    \ifx \bcomment \undefined \def \bcomment#1{#1}\fi
    \ifx \oauthor \undefined \def \oauthor#1{#1}\fi
    \ifx \citeauthoryear \undefined \def \citeauthoryear#1{#1}\fi
    \ifx \endbibitem  \undefined \def \endbibitem {}\fi
    \ifx \bconflocation  \undefined \def \bconflocation#1{#1}\fi
    \ifx \arxivurl  \undefined \def \arxivurl#1{\textsf{#1}}\fi
    \csname PreBibitemsHook\endcsname
    

    \bibitem{alonso2023a}
    \begin{botherref}
    \oauthor{\bsnm{Alonso}, \binits{J.L.}},
    \oauthor{\bsnm{{Bouthelier-Madre}}, \binits{C.}},
    \oauthor{\bsnm{{Clemente-Gallardo}}, \binits{J.}},
    \oauthor{\bsnm{{Mart{\'i}nez-Crespo}}, \binits{D.}}:
    Geometric Flavours of {{Quantum Field}} Theory on a {{Cauchy}} Hypersurface.
      {{Part II}}: {{Canonical}} and {{Geometrical QFT}}
    (2023)
    \end{botherref}
    \endbibitem

    \bibitem{alonso2023d}
    \begin{botherref}
    \oauthor{\bsnm{Alonso}, \binits{J.L.}},
    \oauthor{\bsnm{{Bouthelier-Madre}}, \binits{C.}},
    \oauthor{\bsnm{{Clemente-Gallardo}}, \binits{J.}},
    \oauthor{\bsnm{{Mart{\'i}nez-Crespo}}, \binits{D.}}:
    Hybrid {{Geometrodynamics}}: {{A Hamiltonian}} Description of Classical Gravity
    Coupled to Quantum Matter.
    {arXiv}
    (2023).
    \doiurl{10.48550/arXiv.2307.00922}
    \end{botherref}
    \endbibitem

     %%% 1
     \bibitem{Haag1964}
     \begin{barticle}
     \bauthor{\bsnm{Haag}, \binits{R.}},
     \bauthor{\bsnm{Kastler}, \binits{D.}}:
     \batitle{An {{Algebraic Approach}} to {{Quantum Field Theory}}}.
     \bjtitle{Journal of Mathematical Physics}
     \bvolume{5}(\bissue{7}),
     \bfpage{848}--\blpage{861}
     (\byear{1964}).
     \doiurl{10.1063/1.1704187}
     \end{barticle}
     \endbibitem
     


     %%% 2
     \bibitem{Streater1989}
     \begin{bbook}
     \bauthor{\bsnm{Streater}, \binits{R.F.}},
     \bauthor{\bsnm{Wightman}, \binits{A.S.}}:
     \bbtitle{{{PCT}}, {{Spin}} and {{Statistics}}, and {{All That}}}.
     \bpublisher{{Princeton University Press}}, \blocation{???}
     (\byear{1989}).
     \burl{https://www.jstor.org/stable/j.ctt1cx3vcq}
     Accessed 2023-03-18
     \end{bbook}
     \endbibitem
     
     %%% 3
     \bibitem{Doplicher1969}
     \begin{barticle}
     \bauthor{\bsnm{Doplicher}, \binits{S.}},
     \bauthor{\bsnm{Haag}, \binits{R.}},
     \bauthor{\bsnm{Roberts}, \binits{J.E.}}:
     \batitle{Fields, observables and gauge transformations {{I}}}.
     \bjtitle{Communications in Mathematical Physics}
     \bvolume{13}(\bissue{1}),
     \bfpage{1}--\blpage{23}
     (\byear{1969}).
     \doiurl{10.1007/BF01645267}
     \end{barticle}
     \endbibitem
     
     %%% 4
     \bibitem{Doplicher1969a}
     \begin{barticle}
     \bauthor{\bsnm{Doplicher}, \binits{S.}},
     \bauthor{\bsnm{Haag}, \binits{R.}},
     \bauthor{\bsnm{Roberts}, \binits{J.E.}}:
     \batitle{Fields, observables and gauge transformations. {{II}}}.
     \bjtitle{Communications in Mathematical Physics}
     \bvolume{15}(\bissue{3}),
     \bfpage{173}--\blpage{200}
     (\byear{1969}).
     \doiurl{10.1007/BF01645674}
     \end{barticle}
     \endbibitem
     
     %%% 5
     \bibitem{Doplicher1971}
     \begin{barticle}
     \bauthor{\bsnm{Doplicher}, \binits{S.}},
     \bauthor{\bsnm{Haag}, \binits{R.}},
     \bauthor{\bsnm{Roberts}, \binits{J.E.}}:
     \batitle{Local observables and particle statistics {{I}}}.
     \bjtitle{Communications in Mathematical Physics}
     \bvolume{23}(\bissue{3}),
     \bfpage{199}--\blpage{230}
     (\byear{1971}).
     \doiurl{10.1007/BF01877742}
     \end{barticle}
     \endbibitem
     
     %%% 6
     \bibitem{Doplicher1974}
     \begin{barticle}
     \bauthor{\bsnm{Doplicher}, \binits{S.}},
     \bauthor{\bsnm{Haag}, \binits{R.}},
     \bauthor{\bsnm{Roberts}, \binits{J.E.}}:
     \batitle{Local observables and particle statistics {{II}}}.
     \bjtitle{Communications in Mathematical Physics}
     \bvolume{35}(\bissue{1}),
     \bfpage{49}--\blpage{85}
     (\byear{1974}).
     \doiurl{10.1007/BF01646454}
     \end{barticle}
     \endbibitem
     
     %%% 7
     \bibitem{Buchholz1982}
     \begin{barticle}
     \bauthor{\bsnm{Buchholz}, \binits{D.}},
     \bauthor{\bsnm{Fredenhagen}, \binits{K.}}:
     \batitle{Locality and the structure of particle states}.
     \bjtitle{Communications in Mathematical Physics}
     \bvolume{84}(\bissue{1}),
     \bfpage{1}--\blpage{54}
     (\byear{1982}).
     \doiurl{10.1007/BF01208370}
     \end{barticle}
     \endbibitem
     
     %%% 8
     \bibitem{Buchholz2020}
     \begin{barticle}
     \bauthor{\bsnm{Buchholz}, \binits{D.}},
     \bauthor{\bsnm{Fredenhagen}, \binits{K.}}:
     \batitle{A {{C}}*-algebraic {{Approach}} to {{Interacting Quantum Field
       Theories}}}.
     \bjtitle{Communications in Mathematical Physics}
     \bvolume{377}(\bissue{2}),
     \bfpage{947}--\blpage{969}
     (\byear{2020}).
     \doiurl{10.1007/s00220-020-03700-9}
     \end{barticle}
     \endbibitem
     
     \bibitem{Brunetti2015}
     \begin{bbook}
     \bauthor{\bsnm{Brunetti}, \binits{R.}},
     \bauthor{\bsnm{Dappiaggi}, \binits{C.}},
     \bauthor{\bsnm{Fredenhagen}, \binits{K.}},
     \bauthor{\bsnm{Yngvason}, \binits{J.}}:
     \bbtitle{Advances in {{Algebraic Quantum Field Theory}}}.
     \bpublisher{{Springer International Publishing}},
     \blocation{{Cham}}
     (\byear{2015}).
     \doiurl{10.1007/978-3-319-21353-8}
     \end{bbook}
     \endbibitem

     %%% 9
     \bibitem{connesNoncommutativeGeometry1994}
     \begin{bbook}
     \bauthor{\bsnm{Connes}, \binits{A.}}:
     \bbtitle{Noncommutative {{Geometry}}}.
     \bpublisher{{Academic Press}}
     (\byear{1994})
     \end{bbook}
     \endbibitem
     
     %%% 10
     \bibitem{connesRenormalizationQuantumField2000}
     \begin{barticle}
     \bauthor{\bsnm{Connes}, \binits{A.}},
     \bauthor{\bsnm{Kreimer}, \binits{D.}}:
     \batitle{Renormalization in {{Quantum Field Theory}} and the {{Riemann-Hilbert
       Problem I}}: {{The Hopf Algebra Structure}} of {{Graphs}} and the {{Main
       Theorem}}}.
     \bjtitle{Communications in Mathematical Physics}
     \bvolume{210}(\bissue{1}),
     \bfpage{249}--\blpage{273}
     (\byear{2000}).
     \doiurl{10.1007/s002200050779}
     \end{barticle}
     \endbibitem
     

    %%% 11
    \bibitem{Kibble1979}
    \begin{barticle}
    \bauthor{\bsnm{Kibble}, \binits{T.W.B.}}:
    \batitle{Geometrization of quantum mechanics}.
    \bjtitle{Communications in Mathematical Physics}
    \bvolume{65}(\bissue{2}),
    \bfpage{189}--\blpage{201}
    (\byear{1979}).
    \doiurl{10.1007/BF01225149}
    \end{barticle}
    \endbibitem
    
    %%% 12
    \bibitem{Hida1980}
    \begin{bbook}
    \bauthor{\bsnm{Hida}, \binits{T.}}:
    \bbtitle{Brownian {{Motion}}}.
    \bpublisher{{Springer US}},
    \blocation{{New York, NY}}
    (\byear{1980}).
    \doiurl{10.1007/978-1-4612-6030-1}
    \end{bbook}
    \endbibitem
    
    %%% 13
    \bibitem{Hida1993}
    \begin{bbook}
    \bauthor{\bsnm{Hida}, \binits{T.}},
    \bauthor{\bsnm{Kuo}, \binits{H.-H.}},
    \bauthor{\bsnm{Potthoff}, \binits{J.}},
    \bauthor{\bsnm{Streit}, \binits{L.}}:
    \bbtitle{White {{Noise}}}
    vol. \bseriesno{72}.
    \bpublisher{{Springer Netherlands}},
    \blocation{{Dordrecht}}
    (\byear{1993}).
    \doiurl{10.1007/978-94-017-3680-0}
    \end{bbook}
    \endbibitem
    
    %%% 14
    \bibitem{Wester2003}
    \begin{botherref}
    \oauthor{\bsnm{Westerkamp}, \binits{W.}}:
    Recent {{Results}} in {{Infinite Dimensional Analysis}} and {{Applications}} to
      {{Feynman Integrals}}
    (2003)
    {\href{https://arxiv.org/abs/math-ph/0302066}{{arXiv:math-ph/0302066}}}
    \end{botherref}
    \endbibitem
    

    %%% 16
    \bibitem{Holden2010}
    \begin{bbook}
    \bauthor{\bsnm{Holden}, \binits{H.}},
    \bauthor{\bsnm{{\O}ksendal}, \binits{B.}},
    \bauthor{\bsnm{Ub{\o}e}, \binits{J.}},
    \bauthor{\bsnm{Zhang}, \binits{T.}}:
    \bbtitle{Stochastic {{Partial Differential Equations}}}
    vol. \bseriesno{7}.
    \bpublisher{{Springer New York}},
    \blocation{{New York, NY}}
    (\byear{2010}).
    \doiurl{10.1007/978-0-387-89488-1}
    \end{bbook}
    \endbibitem

    %%% 17
    \bibitem{Kondratiev1996}
    \begin{barticle}
    \bauthor{\bsnm{Kondratiev}, \binits{Y.G.G.}},
    \bauthor{\bsnm{Leukert}, \binits{P.}},
    \bauthor{\bsnm{Potthoff}, \binits{J.}},
    \bauthor{\bsnm{Streit}, \binits{L.}},
    \bauthor{\bsnm{Westerkamp}, \binits{W.}}:
    \batitle{Generalized {{Functionals}} in {{Gaussian Spaces}}: {{The
    Characterization Theorem Revisited}}}.
    \bjtitle{Journal of Functional Analysis}
    \bvolume{141}(\bissue{2}),
    \bfpage{301}--\blpage{318}
    (\byear{1996})
    {\href{https://arxiv.org/abs/math/0303054}{{arXiv:math/0303054}}}.
    \doiurl{10.1006/jfan.1996.0130}
    \end{barticle}
    \endbibitem

    %%% 15
    \bibitem{Nunno2009}
    \begin{bbook}
    \bauthor{\bsnm{Nunno}, \binits{G.D.}},
    \bauthor{\bsnm{{\O}ksendal}, \binits{B.}},
    \bauthor{\bsnm{Proske}, \binits{F.}}:
    \bbtitle{Malliavin {{Calculus}} for {{L\'evy Processes}} with {{Applications}}
        to {{Finance}}}
    vol. \bseriesno{44}.
    \bpublisher{{Springer Berlin Heidelberg}},
    \blocation{{Berlin, Heidelberg}}
    (\byear{2009}).
    \doiurl{10.1007/978-3-540-78572-9}
    \end{bbook}
    \endbibitem

    %%% 16
    \bibitem{Corichi2004}
    \begin{barticle}
    \bauthor{\bsnm{Corichi}, \binits{A.}},
    \bauthor{\bsnm{Cortez}, \binits{J.}},
    \bauthor{\bsnm{Quevedo}, \binits{H.}}:
    \batitle{Schr\"odinger and {{Fock}} representation for a field theory on curved
      spacetime}.
    \bjtitle{Annals of Physics}
    \bvolume{313}(\bissue{2}),
    \bfpage{446}--\blpage{478}
    (\byear{2004}).
    \doiurl{10.1016/j.aop.2004.05.004}
    \end{barticle}
    \endbibitem
    

    
    %%% 18
    \bibitem{Oeckl2012a}
    \begin{barticle}
    \bauthor{\bsnm{Oeckl}, \binits{R.}}:
    \batitle{Holomorphic quantization of linear field theory in the general
      boundary formulation}.
    \bjtitle{Symmetry, Integrability and Geometry: Methods and Applications
      (SIGMA)}
    \bvolume{8},
    \bfpage{31}
    (\byear{2012})
    {\href{https://arxiv.org/abs/1009.5615}{{arXiv:1009.5615}}}.
    \doiurl{10.3842/SIGMA.2012.050}
    \end{barticle}
    \endbibitem
    
    %%% 17
    \bibitem{Oeckl2012}
    \begin{barticle}
    \bauthor{\bsnm{Oeckl}, \binits{R.}}:
    \batitle{The {{Schr\"odinger}} representation and its relation to the
        holomorphic representation in linear and affine field theory}.
    \bjtitle{Journal of Mathematical Physics}
    \bvolume{53}(\bissue{7}),
    \bfpage{072301}
    (\byear{2012})
    {\href{https://arxiv.org/abs/1109.5215}{{arXiv:1109.5215}}}.
    \doiurl{10.1063/1.4731770}
    \end{barticle}
    \endbibitem

    %18.5
    \bibitem{henry-labordere2009}
    \begin{bbook}
    \bauthor{\bsnm{{Henry-Labord{\`e}re}}, \binits{P.}}:
    \bbtitle{Analysis, Geometry, and Modeling in Finance: Advanced Methods in
    Option Pricing}.
    \bpublisher{{CRC Press}},
    \blocation{{Boca Raton}}
    (\byear{2009})
    \end{bbook}
    \endbibitem
    %%% 19
    \bibitem{Gellfand}
    \begin{botherref}
    \oauthor{\bsnm{Moiseevi{\v c}~Gel'fand}, \binits{I.}},
    \oauthor{\bsnm{Vilenkin}, \binits{N.Y.}}:
    Generalized functions: {{Applications}} of {{Harmonic Analysis}}, {{Volume}} 4
    (1967)
    \end{botherref}
    \endbibitem
    

    %%% 19
    \bibitem{dito1990}
    \begin{barticle}
    \bauthor{\bsnm{Dito}, \binits{J.}}:
    \batitle{Star-product approach to quantum field theory: {{The}} free scalar
    field}.
    \bjtitle{Letters in Mathematical Physics}
    \bvolume{20}(\bissue{2}),
    \bfpage{125}--\blpage{134}
    (\byear{1990}).
    \doiurl{10.1007/BF00398277}
    \end{barticle}
    \endbibitem
    %%% 20
    \bibitem{Kriegl1997}
    \begin{bbook}
    \bauthor{\bsnm{Kriegl}, \binits{A.}},
    \bauthor{\bsnm{Michor}, \binits{P.}}:
    \bbtitle{The {{Convenient Setting}} of {{Global Analysis}}}
    vol. \bseriesno{53}.
    \bpublisher{{American Mathematical Society}},
    \blocation{{Providence, Rhode Island}}
    (\byear{1997}).
    \doiurl{10.1090/surv/053}
    \end{bbook}
    \endbibitem
    
    %%% 21
    \bibitem{gotayMomentumMapsClassical2004a}
    \begin{botherref}
    \oauthor{\bsnm{Gotay}, \binits{M. J.}},
    \oauthor{\bsnm{Isenberg}, \binits{ J.}},
    \oauthor{\bsnm{Marsden}, \binits{ J. E.}},
    \oauthor{\bsnm{Montgomery}, \binits{ R}}:
    Momentum {{Maps}} and {{Classical Relativistic Fields}}. {{Part I}}: {{Covariant Field Theory}}
    (2004)
    {\href{https://arxiv.org/abs/physics/9801019}{{arXiv:physics/9801019}}}
    \end{botherref}
    \endbibitem
    
       %%% 22
       \bibitem{gotayMomentumMapsClassical2004}
       \begin{botherref}
       \oauthor{\bsnm{Gotay}, \binits{M. J.}},
       \oauthor{\bsnm{Isenberg}, \binits{ J.}},
       \oauthor{\bsnm{Marsden}, \binits{ J. E.}},
       \oauthor{\bsnm{Montgomery}, \binits{ R}}:
       Momentum {{Maps}} and {{Classical Relativistic Fields}}. {{Part II}}: {{Canonical Analysis}} of {{Field Theories}}
       (2004)
       {\href{https://arxiv.org/abs/math-ph/0411032}{{arXiv:math-ph/0411032}}}
       \end{botherref}
       \endbibitem
    \endbibitem


    %%% 23.5
    \bibitem{Moshayedi2019}
    \begin{botherref}
    \oauthor{\bsnm{Moshayedi}, \binits{N.}}:
    Quantum field theory and functional integrals.
    arXiv
    (2019)
    {\href{https://arxiv.org/abs/1902.08652}{{arxiv:1902.08652}}}
    \end{botherref}
    \endbibitem

    %%% 23

    

    
    %%% 25
    \bibitem{Woodhouse1997}
    \begin{bbook}
    \bauthor{\bsnm{Woodhouse}, \binits{N.M.J.}}:
    \bbtitle{Geometric {{Quantization}}},
    \bedition{Second edition} edn.
    \bsertitle{Oxford {{Mathematical Monographs}}}.
    \bpublisher{{Oxford University Press}},
    \blocation{{Oxford}}
    (\byear{1997})
    \end{bbook}
    \endbibitem
    
    %%% 26
    \bibitem{hall2013}
    \begin{bbook}
    \bauthor{\bsnm{Hall}, \binits{B.C.}}:
    \bbtitle{Quantum Theory for Mathematicians}.
    \bsertitle{Graduate Texts in Mathematics},
    vol. \bseriesno{267}.
    \bpublisher{{Springer}},
    \blocation{{New York}}
    (\byear{2013})
    \end{bbook}
    \endbibitem
    
    \end{thebibliography}


    

    \end{document}
